\documentclass[11pt]{article}
\usepackage[T1]{fontenc}
\usepackage[utf8]{inputenc}
\usepackage{lmodern,microtype}
\usepackage[a4paper,top=22mm,bottom=23mm,left=23mm,right=23mm]{geometry}
\usepackage{amsmath,amssymb,amsthm,mathtools,mathrsfs}
\usepackage{booktabs,longtable,array,tabularx,multirow}
\usepackage[table]{xcolor}
\usepackage{graphicx,tikz}
\usetikzlibrary{arrows.meta,positioning,calc,fit}
\usepackage{needspace,fancyhdr,titlesec,enumitem,float}
\usepackage[most]{tcolorbox}
\usepackage{hyperref,xurl}
\definecolor{Navy}{HTML}{244F7B}
\definecolor{Teal}{HTML}{167C80}
\definecolor{Amber}{HTML}{B77623}
\definecolor{Ice}{HTML}{F3F6FA}
\definecolor{Ink}{HTML}{202B36}
\hypersetup{colorlinks=true,linkcolor=Navy,citecolor=Teal,urlcolor=Navy,
  pdftitle={Sharp Gaussian Khatri--Rao Embeddings: Three Proof Chains: Concentration, Framework, and Wick--Hermite},
  pdfauthor={Method Problems -- consolidated research note},bookmarksnumbered=true}
\color{Ink}
\pagestyle{fancy}\fancyhf{}
\fancyhead[L]{\sffamily\small Gaussian Khatri--Rao embeddings}
\fancyhead[R]{\sffamily\small Three proof chains: A / B / C}
\fancyfoot[L]{\sffamily\footnotesize Three-chain comparison\;\textbullet\;16 September 2026}
\fancyfoot[R]{\sffamily\small\thepage}
\renewcommand{\headrulewidth}{.4pt}
\setlength{\headheight}{15pt}
\setlength{\emergencystretch}{2em}
\setlength{\parskip}{3pt}
\setlength{\parindent}{0pt}
\linespread{1.04}
\allowdisplaybreaks[1]
\clubpenalty=10000\widowpenalty=10000
\titleformat{\section}{\Large\sffamily\bfseries\color{Navy}}{\thesection}{.65em}{}
\titleformat{\subsection}{\large\sffamily\bfseries\color{Navy}}{\thesubsection}{.6em}{}
\titleformat{\subsubsection}{\normalsize\sffamily\bfseries}{\thesubsubsection}{.6em}{}
\titlespacing*{\section}{0pt}{2.0ex plus 1ex minus .2ex}{1.0ex}
\titlespacing*{\subsection}{0pt}{1.8ex plus 1ex minus .2ex}{.7ex}
\numberwithin{equation}{section}
\newtheorem{theorem}{Theorem}[section]
\newtheorem{lemma}[theorem]{Lemma}
\newtheorem{proposition}[theorem]{Proposition}
\newtheorem{corollary}[theorem]{Corollary}
\theoremstyle{definition}\newtheorem{definition}[theorem]{Definition}
\theoremstyle{remark}\newtheorem{remark}[theorem]{Remark}
\newcommand{\R}{\mathbb R}\newcommand{\E}{\mathbb E}\newcommand{\Prob}{\mathbb P}
\newcommand{\tr}{\operatorname{tr}}\newcommand{\Ent}{\operatorname{Ent}}
\newcommand{\norm}[1]{\lVert#1\rVert}\newcommand{\ip}[2]{\langle#1,#2\rangle}
\newcommand{\one}{\mathbf1}\newcommand{\F}{\mathrm F}
\newcommand{\step}[1]{\par\addvspace{.65\baselineskip}\Needspace{4\baselineskip}\noindent\textbf{#1}\par\nobreak\smallskip}
\newcolumntype{Y}{>{\raggedright\arraybackslash}X}
\newcolumntype{P}[1]{>{\raggedright\arraybackslash}p{#1}}
\newtcolorbox{takeaway}[1][]{colback=Ice,colframe=Navy!50,boxrule=.45pt,
  arc=1pt,left=8pt,right=8pt,top=7pt,bottom=7pt,fonttitle=\sffamily\bfseries,
  coltitle=Navy,colbacktitle=Ice,enhanced,#1}
\setlist{nosep,leftmargin=1.5em}
\DeclareMathOperator{\diag}{diag}
\DeclareFontFamily{U}{rsfs}{}
\DeclareFontShape{U}{rsfs}{m}{n}{<-6> rsfs5 <6-8> rsfs7 <8-> rsfs10}{}
\begin{document}
\thispagestyle{empty}
{\sffamily\color{Navy}\small METHOD PROBLEMS\hfill PROOF-CHAIN MAP\par}
\vspace{3mm}
{\sffamily\bfseries\fontsize{24}{28}\selectfont Sharp Gaussian Khatri--Rao embeddings\par}
\vspace{1.5mm}
{\sffamily\large Three proof chains: where they split and meet\par}
\vspace{2mm}
{\small Read downward. A and B retain the original concentration and framework proofs.
C is the separate Wick--Hermite hybrid. The model and final sample-size order
are shared; the intermediate estimates and constants are not.\par}
\vspace{1mm}
\newlength{\KRFullWidth}\setlength{\KRFullWidth}{\dimexpr\textwidth-16pt\relax}
\newlength{\KRColumnWidth}\setlength{\KRColumnWidth}{.333333\textwidth}
\addtolength{\KRColumnWidth}{-4mm}\addtolength{\KRColumnWidth}{-12pt}
\newlength{\KRColumnOffset}\setlength{\KRColumnOffset}{.333333\textwidth}
\addtolength{\KRColumnOffset}{2mm}
\begin{center}
% Physical-size vector artwork: no rasterization, scaling, or generated image.
\begin{tikzpicture}[
 >={Latex[length=1.6mm,width=1.2mm]},
 every node/.style={font=\fontsize{10}{12}\selectfont},
 base/.style={rounded corners=1.5pt,line width=.5pt,align=center,
  inner xsep=5pt,inner ysep=4pt,outer sep=.5pt},
 shared/.style={base,draw=Navy!70,fill=Ice,text width=\KRFullWidth},
 abox/.style={base,draw=Teal!85,fill=Teal!4,text width=\KRColumnWidth,minimum height=19mm},
 bbox/.style={abox,draw=Navy!75,fill=Navy!3},
 cbox/.style={abox,draw=Amber!85,fill=Amber!5},
 ahead/.style={base,draw=Teal,fill=Teal,text=white,text width=\KRColumnWidth,
  minimum height=11mm,font=\sffamily\bfseries\fontsize{10}{12}\selectfont},
 bhead/.style={ahead,draw=Navy,fill=Navy},
 chead/.style={ahead,draw=Amber,fill=Amber},
 scope/.style={shared,draw=Amber!80,fill=Amber!7},
 flow/.style={->,draw=Navy!85,line width=.6pt},
 aflow/.style={flow,draw=Teal!90},cflow/.style={flow,draw=Amber!95}]
\node[shared,minimum height=23mm] (geometry)
 {\textbf{Shared input: the original Gaussian law and arbitrary fixed isometry $U$}\\[2pt]
  $y_i=U^{\mathsf T}(g_i\otimes h_i),\quad Y_i=y_iy_i^{\mathsf T},
  \quad Z=m^{-1}\sum_i(Y_i-I_r)$\\[2pt]
  $A(z)w=U^{\mathsf T}(z\otimes w),\quad \|A(z)\|\le\|z\|,
  \quad\mathbb E\|A(g)\|_F^2=r$};
\node[ahead,below=7mm of geometry.south,xshift=-\KRColumnOffset] (ah)
 {CHAIN A\\Matrix concentration};
\node[bhead,below=7mm of geometry.south] (bh)
 {CHAIN B\\Original framework};
\node[chead,below=7mm of geometry.south,xshift=\KRColumnOffset] (ch)
 {CHAIN C\\Wick--Hermite hybrid};
\coordinate (split) at ($(geometry.south)+(0,-3mm)$);
\draw[aflow] (geometry.south)--(split)-|(ah.north);
\draw[flow] (geometry.south)--(bh.north);
\draw[cflow] (geometry.south)--(split)-|(ch.north);

\node[abox,below=3mm of ah] (a1)
 {\textbf{A1. Matrix powers}\\[2pt]
  $Y^p=\|y\|^{2p-2}yy^{\mathsf T}$\\
  Retain matrix orientation};
\node[bbox,below=3mm of bh] (b1)
 {\textbf{B1. Original-law lift}\\[2pt]
  $\mathbb E\tr Z^{2q}$: vacuum energy\\
  Multiply by the original $Z$};
\node[cbox,below=3mm of ch] (c1)
 {\textbf{C1. Wick radial moments}\\[2pt]
  $\mathbb E\|y\|^{2p}\le\kappa_pD_{r,p}$\\
  Joint Gaussian integration};
\draw[aflow] (ah)--(a1);\draw[flow] (bh)--(b1);\draw[cflow] (ch)--(c1);

\node[abox,below=3mm of a1] (a2)
 {\textbf{A2. Auxiliary truncation}\\[2pt]
  $R_s=72(r+s)$\\
  Charge clipping and bias};
\node[bbox,below=3mm of b1] (b2)
 {\textbf{B2. Exact cutoff}\\[2pt]
  Hermite degree $\le4q$; $k\le q$\\
  Keep even factor symmetries};
\node[cbox,below=3mm of c1] (c2)
 {\textbf{C2. Mixed contractions}\\[2pt]
  $\mathbb EX^2\preceq(4r+4)I_r$\\
  Keep coefficient orthogonality};
\draw[aflow] (a1)--(a2);\draw[flow] (b1)--(b2);\draw[cflow] (c1)--(c2);

\node[abox,below=3mm of a2] (a3)
 {\textbf{A3. Matrix Laplace}\\[2pt]
  Lieb--Tropp + mean offset\\
  Raw proxy $9r$; scale $R_s$};
\node[bbox,below=3mm of b2] (b3)
 {\textbf{B3. Baseline contractions}\\[2pt]
  Fresh: $9r-1$; repeat: $B_q$\\
  Sum local degrees first};
\node[cbox,below=3mm of c2] (c3)
 {\textbf{C3. Wick into Hermite}\\[2pt]
  Same exact cutoff as B2\\
  Sharpen repeat cost $\widetilde B_q$};
\draw[aflow] (a2)--(a3);\draw[flow] (b2)--(b3);\draw[cflow] (c2)--(c3);

\node[abox,below=3mm of a3] (a4)
 {\textbf{A4. Original-law tail}\\[2pt]
  $U_A(m,r,\delta)\le\varepsilon$\\
  Original, unmodified sketch};
\node[bbox,below=3mm of b3] (b4)
 {\textbf{B4. Vacuum + Markov}\\[2pt]
  $r\|G_B^{[q]}\cdots G_B^{[1]}e_0\|^2$\\
  Keep stage order and vacuum};
\node[cbox,below=3mm of c3] (c4)
 {\textbf{C4. Refined vacuum}\\[2pt]
  $r\|\widehat G_C^{[q]}\cdots\widehat G_C^{[1]}e_0\|^2$\\
  Exact radial + degree costs};
\draw[aflow] (a3)--(a4);\draw[flow] (b3)--(b4);\draw[cflow] (c3)--(c4);

\node[fit=(a4)(b4)(c4),inner sep=0pt,outer sep=0pt] (lastrow) {};
\node[shared,below=6mm of lastrow.south,minimum height=18mm] (upper)
 {\textbf{Convergence: one upper-bound order, three distinct certificates}\\[3pt]
 $m\ge C\bigl[\varepsilon^{-2}rt+\varepsilon^{-1}t^2\bigr]
 \ \Longrightarrow\ \mathbb P\{\|Z\|>\varepsilon\}\le\delta,
 \qquad t=\log(4r/\delta)$};
\draw[aflow] (a4.south)--++(0,-2.6mm)-|([xshift=-\KRColumnOffset]upper.north);
\draw[flow] (b4.south)--(upper.north);
\draw[cflow] (c4.south)--++(0,-2.6mm)-|([xshift=\KRColumnOffset]upper.north);
\node[scope,below=3.5mm of upper,minimum height=18mm] (lower)
 {\textbf{Combine with the separately proved, common lower bound}\\[2pt]
  $y_i=\eta_i\xi_i$ forces $m=\Omega(r\log r)$ in the worst case.\\
  All three chains attain $\Theta(r\log(er))$ at fixed distortion and confidence.};
\draw[flow,dashed,draw=Amber] (upper)--(lower);
\end{tikzpicture}
\end{center}
\vfill
{\footnotesize\color{Ink!75}Here $\kappa_p=(2p-1)!!$ and
$D_{r,p}=\prod_{\ell=0}^{p-1}(r+2\ell)$. Page 2 fixes the A/B/C comparison.
The Wick-assisted concentration variant is documented separately, not substituted
into A or labelled C. All application-specific proofs follow in this document.}
\clearpage
\section*{Three-chain comparison}
\phantomsection\label{sec:three-chain-overview}
{\sffamily\color{Navy}\large Same Khatri--Rao problem; three visibly separate proof routes.\par}
\vspace{2mm}
The main labels are fixed throughout: \textbf{A = original matrix concentration},
\textbf{B = original framework}, and \textbf{C = combined Wick--Hermite framework}.
C has its own local estimates, proof section, stage matrices, and numerical certificate.

\begingroup\fontsize{9.4}{11.4}\selectfont
\setlength{\tabcolsep}{4pt}\renewcommand{\arraystretch}{1.13}
\begin{tabularx}{\textwidth}{@{}P{.14\textwidth}YYY@{}}
\toprule
\textbf{Comparison}&{\color{Teal}\textbf{Chain A} (\S\ref{sec:A})}&
{\color{Navy}\textbf{Chain B} (\S\ref{sec:B})}&
{\color{Amber}\textbf{Chain C} (\S\ref{sec:C})}\\
\midrule
\rowcolor{Ice}
\textbf{Retained object}&One-row matrix powers $\mathbb E Y^p$.&
Original-law vacuum walks, sample supports, and local degrees.&
Exact Wick radial/mixed moments inserted into the same Hermite walks.\\[3pt]
\textbf{Domain}&An auxiliary rare-outcome truncation; restore the original law.&
Exact cutoff: total degree $\le4q$, active count $\le q$, even symmetries.&
The same exact cutoff. No outcome truncation and no independent replacement.\\[3pt]
\rowcolor{Ice}
\textbf{Gaussian input}&Collective Lipschitz row estimate and directional fourth moments.&
The same baseline row estimate; H\"older and hypercontractivity at local degree $d$.&
Sharp radial product $\kappa_pD_{r,p}$ and orthogonal mixed-fourth contractions.\\[3pt]
\textbf{Variance used}&Raw second-moment proxy $9r$.&
Centered activation proxy $9r-1$.&
Centered activation proxy $4r+4$.\\[3pt]
\rowcolor{Ice}
\textbf{Assembly}&Lieb--Tropp matrix Laplace with a deterministic mean offset.&
Positive sector majorant and ordered vacuum propagation, then Markov.&
The same propagation, with smaller costs and refined active-degree budgets.\\[3pt]
\textbf{Explicit output}&Tail allowance $U_A$ in \eqref{eq:A-tail}.&
Moment and tail certificates in \eqref{eq:stage-moment}.&
Refined moment certificate \eqref{c:eq:wick-stage-moment}.\\[3pt]
\rowcolor{Ice}
\textbf{What is not claimed}&An optimal concentration constant.&
Necessity of the framework for this problem.&
A new sketch, a new sample-size order, or a pure Wick-only proof.\\[3pt]
\bottomrule
\end{tabularx}
\endgroup

\vspace{3mm}
{\sffamily\bfseries\color{Navy}Explicit sufficient row counts: A versus B versus C\par}
{\small $\varepsilon=1/2$, $\delta=0.01$; identical model and subspace scope.
B uses its original stage-dependent certificate; C uses the radial/degree refinement.}
\begin{table}[H]\centering\small
\setlength{\tabcolsep}{10pt}\renewcommand{\arraystretch}{1.18}
\begin{tabular}{@{}rrrrr@{}}
\toprule
$r$&\textbf{A: concentration}&\textbf{B: framework}&\textbf{C: Wick--Hermite}&$q_B=q_C$\\
\midrule
10&91\,000&60\,000&12\,800&3\\
100&580\,000&472\,000&111\,000&5\\
1\,000&6\,200\,000&5\,150\,000&1\,250\,000&7\\
10\,000&72\,000\,000&61\,000\,000&15\,000\,000&8\\
\bottomrule
\end{tabular}
\caption{The three main certificates. These are selected sufficient theorem bounds,
not empirical row counts or global minima. Parameters and interval checks appear in
Section~\ref{sec:comparison} and Appendix~\ref{app:checks}.}
\label{tab:three-main-counts}
\end{table}

\begin{takeaway}[title={What this comparison means},top=5pt,bottom=5pt]
\small
All three prove the same rate. The change from B to C measures the Wick--Hermite
refinement while keeping the final moment order fixed. It does not establish
universal superiority over concentration: the separate Wick-assisted variant
$A^{\mathrm W}$ also improves A, and is compared with C in
Appendix~\ref{app:A-wick}. It is not a fourth main chain.
\end{takeaway}
\clearpage
\section{Model, theorem, and comparison at a glance}\label{sec:model}
\subsection{The unchanged sketch and its quantifiers}
Let $n_1,n_2,r,m\ge1$ be integers with $r\le n_1n_2$, and fix a deterministic isometry
\[
 U:\R^r\longrightarrow\R^{n_1}\otimes\R^{n_2},\qquad U^{\mathsf T}U=I_r.
\]
For $1\le i\le m$, let $g_i\sim N(0,I_{n_1})$ and $h_i\sim N(0,I_{n_2})$;
all samples and both factors are mutually independent. Define
\begin{equation}\label{eq:model}
 \begin{aligned}
 S&=\frac1{\sqrt m}\begin{pmatrix}(g_1\otimes h_1)^{\mathsf T}\\\vdots\\(g_m\otimes h_m)^{\mathsf T}\end{pmatrix},
 &y_i&=U^{\mathsf T}(g_i\otimes h_i),\\
 Y_i&=y_iy_i^{\mathsf T},\qquad X_i=Y_i-I_r,
 &Z&=U^{\mathsf T}S^{\mathsf T}SU-I_r=\frac1m\sum_{i=1}^mX_i.
 \end{aligned}
\end{equation}
Independence gives $\E(g_i\otimes h_i)(g_i\otimes h_i)^{\mathsf T}=I_{n_1n_2}$,
so $\E Y_i=I_r$ and $\E X_i=0$.
The event $\norm Z\le\varepsilon$ is exactly the squared-norm embedding guarantee
\[
 (1-\varepsilon)\norm x^2\le\norm{SUx}^2\le(1+\varepsilon)\norm x^2
 \quad\text{for every }x\in\R^r.
\]
Every probability statement below is for \emph{each fixed} $U$, with constants uniform
in $U,n_1,n_2$. It is not a simultaneous statement for all subspaces after $S$ is drawn.
All logarithms are natural; unmarked matrix norms are operator norms.

\begin{theorem}[Common endpoint]\label{thm:common-endpoint}
There is a universal $C>0$ such that, for $0<\varepsilon,\delta<1$ and
$t=\log(4r/\delta)$,
\begin{equation}\label{eq:main-size}
 m\ge C\left(\frac{rt}{\varepsilon^2}+\frac{t^2}{\varepsilon}\right)
 \quad\Longrightarrow\quad\Prob\{\norm Z>\varepsilon\}\le\delta.
\end{equation}
Each of Chains A, B, and C below establishes this upper bound. A and B use the
baseline local estimates; C replaces them by its Wick contractions and then
uses the common original-law Hermite interface. Conversely, for $Ue_a=e_a\otimes e_1$ with $n_1=r,n_2=2$, there is
$c>0$ such that uniformly over integers $1\le m\le c r\log r$,
\[
 \Prob\{\norm Z>1/2\}\longrightarrow1\qquad(r\longrightarrow\infty).
\]
Therefore the worst-case sample order at distortion $1/2$ and success probability
$0.99$ is $\Theta(r\log(er))$.
\end{theorem}
The statement does not assert optimal joint dependence on $\varepsilon$ and $\delta$.
The upper bound is proved for arbitrary subspace coefficients; only the lower-bound
example uses a common factor.

\subsection{Reading guide and chain conventions}
Section~\ref{sec:common} fixes the original local estimates for A and B.
Sections~\ref{sec:A}, \ref{sec:B}, and \ref{sec:C} are the three proof chains.
C reuses B's exact multiplication/cutoff identities, but supplies its own Wick
contractions, repeated-visit budgets, and refined propagation theorem. It is
not a renaming of B. The lower bound is shared in Section~\ref{sec:lower}.

Section~\ref{sec:comparison} compares fixed A/B/C certificates. The sharper
concentration variant $A^{\mathrm W}$ is separated in Appendix~\ref{app:A-wick}.
Classical Gaussian inputs are derived in Appendix~\ref{app:gaussian}; Lieb's
trace concavity is stated and attributed in Section~\ref{sec:A-laplace}.
No earlier project note is needed for an application-specific proof step.

\clearpage
\subsection{Related work and attribution}\label{sec:attribution-intro}
The sketch is not a new construction. Tensor random projections were studied by
Sun, Guo, Tropp, and Udell~\cite{Sun}, and fixed-vector polynomial-kernel sketching
bounds appear in Ahle et al.~\cite[Appendix~A.2, Theorem~42]{Ahle}.
In particular, the order of the $r=1$ sample bound is already present in that
fixed-vector literature; the target
here is uniformity over all directions of an arbitrary fixed $r$-dimensional
subspace. Bujanovi\'c and Kressner~\cite{BK} study norm and trace estimation with
random rank-one vectors. Bujanovi\'c, Grubi\v{s}i\'c, Kressner, and
Lam~\cite{BGKL} prove an earlier OSE theorem for this two-factor Gaussian model.

Beretta--Musco~\cite[Theorem~1.3]{BM} give
$O(\varepsilon^{-2}r\log^3(r/(\varepsilon\delta)))$ rows in the Gaussian
order-two specialization. Their theorem covers more distributions and tensor
orders. The comparison here is restricted to that specialization, with an
absolute rescaling of distortion when passing between norm and squared-norm
conventions. Neither the present constants nor the numerical table are advertised
as universal improvements over their broader result.

\paragraph{A version-sensitive predecessor.}
Saibaba, Verma, and Ballard's first version~\cite[Theorem~11]{SVBv1}
contains a near-linear, polylogarithmic OSE claim and an accompanying proof
attempt. The second version~\cite{SVBv2} omits that assertion. Beretta--Musco's
acknowledgments~\cite{BM} report that a proof error prompted its removal and
credit the earlier strategy. We retain this historical credit but do not use the
removed assertion as a theorem. The paper as a whole is not described as withdrawn.

\paragraph{Methodological antecedents.}
Chain A follows the established truncation-and-concentration architecture for
Khatri--Rao sketches; compare~\cite[Section~2.3]{BM}. Its specific calculations
retain one-row matrix moments and the distinction between the truncation height
and the exponential-moment scale. The matrix Laplace mechanism is the
Lieb--Tropp method, not a new concentration theorem~\cite{Lieb,Tropp}.
The order of the shared projected-row estimate also follows from the Hilbert-space
specialization of Adamczak--Lata\l a--Meller~\cite[Corollary~12]{ALM};
Section~\ref{sec:common} supplies an explicit normalization and a short direct proof.
Coefficient-sensitive matrix-chaos estimates and Khatri--Rao applications also
precede this note in Bandeira, Lucca, Nizi\'c-Nikolac, and
van Handel~\cite{BLNvH}. Their full-matrix norm estimates are not identified with
the centered, arbitrary-fixed-subspace covariance theorem proved here.

Chain B is organized using the supplied framework~\cite{Framework}. Its
multiplication representation, Hermite decomposition, and Jacobi moment
interpretation have classical antecedents~\cite{Wiener,AccardiBozejko,GolubWelsch}.
The application-specific content to assess is the sequence of local estimates,
the resulting sample bound, and its explicit vacuum certificates, rather than
ownership of these representation principles. Section~\ref{sec:provenance-table}
separates the sources and roles. No publication-priority claim is made.

\paragraph{The hybrid is separated from both baselines.}
Chain C combines the classical Wick--Isserlis pairing rule with the Hermite
and energy interface of B~\cite{Isserlis,Janson,Framework}. The radialization
mechanism also appears in the supplied attributed rMPS comparison~\cite{RMPS};
that manuscript is cited for internal provenance, not as a theorem about the
present Khatri--Rao model. The application-specific refinements were recorded
in the supplied Wick supplement~\cite{WickNote} and are proved here in
Section~\ref{sec:C}. The rMPS comparison is used as a layout reference only:
its reconstruction model, inverse-Gram expansions, and numerical tables are not
transferred to the present theorem.
\clearpage
\section{Shared baseline coefficient estimates}\label{sec:common}
\begin{takeaway}
The same subspace geometry feeds all three chains. This section fixes the
pre-Wick estimates used by A and B; C supplies sharper local inputs separately. Keep the map $A(z)$ collected before
taking norms: its operator-norm Lipschitz constant is $1$, not $\sqrt r$.
Then keep the rank-one orientation when a matrix-valued estimate is needed.
\end{takeaway}

\subsection{Gaussian moment control with explicit normalization}
Write $\gamma_N$ for standard Gaussian measure on $\R^N$ and
$\Ent(F)=\E[F\log F]-(\E F)\log\E F$ for $F\ge0$.
The Gaussian logarithmic Sobolev inequality is
\begin{equation}\label{eq:LSI}
 \Ent_{\gamma_N}(f^2)\le2\E_{\gamma_N}\norm{\nabla f}^2.
\end{equation}
For a Hilbert-valued homogeneous Hermite polynomial $F_d$ of total degree $d$,
Gaussian hypercontractivity gives
\begin{equation}\label{eq:HC}
 \norm{F_d}_{L_u}\le(u-1)^{d/2}\norm{F_d}_{L_2},\qquad u\ge2.
\end{equation}
The normalization and short semigroup proofs appear in Appendix~\ref{app:gaussian};
these are classical Nelson--Gross inequalities~\cite{Nelson,Gross,GrossSurvey}.
Differentiating moments under a logarithmic Sobolev inequality is also a standard
route to moment estimates; see Aida--Stroock~\cite{AidaStroock}.
The following lemma fixes the Gaussian normalization used throughout.

\begin{lemma}[Lipschitz moments]\label{lem:lip}
For a nonnegative $L$-Lipschitz function $f$ on a finite-dimensional standard
Gaussian space and any real $p\ge2$,
\begin{equation}\label{eq:lip}
 \norm f_{L_p}^2\le\norm f_{L_2}^2+(p-2)L^2.
\end{equation}
\end{lemma}
\begin{proof}
Assume first that $f$ is smooth, bounded, and bounded below by a positive
constant, and set $M(p)=\E f^p$.

\step{1. Differentiate the running moment in its exponent.}
Since $\frac{d}{dp}\E f^p=\E f^p\log f$, a direct computation gives
\[
 \frac{d}{dp}M(p)^{2/p}
 =\frac2{p^2}M(p)^{2/p-1}\Ent(f^p).
\]

\step{2. Bound the entropy by the logarithmic Sobolev inequality.}
Apply \eqref{eq:LSI} to $f^{p/2}$ and use $\norm{\nabla f}\le L$:
\[
 \Ent(f^p)\le\frac{p^2L^2}2\,\E f^{p-2}.
\]
H\"older's inequality gives $\E f^{p-2}\le M(p)^{(p-2)/p}$, so the two previous
displays combine into
\[
 \frac{d}{dp}M(p)^{2/p}
 \le L^2M(p)^{2/p-1}M(p)^{1-2/p}=L^2.
\]

\step{3. Integrate from $2$ to $p$.}
This yields $M(p)^{2/p}\le M(2)+(p-2)L^2$, which is \eqref{eq:lip}.
Smoothing, a positive regularization, and truncation of the scalar test function
justify the same bound for Lipschitz $f$ by passage to the limit.
\end{proof}

\subsection{The collective conditional Gaussian map}
Reshape the columns of $U$ as matrices $B_1,\ldots,B_r\in\R^{n_1\times n_2}$.
Then
\begin{equation}\label{eq:B-orthonormal}
 \tr(B_a^{\mathsf T}B_b)=\delta_{ab},\qquad
 y_a=g^{\mathsf T}B_a h.
\end{equation}
For $z\in\R^{n_1}$ define $A(z):\R^{n_2}\to\R^r$ by
\begin{equation}\label{eq:A-map}
 A(z)w=U^{\mathsf T}(z\otimes w),\qquad y=A(g)h.
\end{equation}
Because $U^{\mathsf T}$ is a contraction,
\begin{equation}\label{eq:collective}
 \norm{A(z)}\le\norm z,\qquad
 \E\norm{A(g)}_\F^2=r,\qquad
 \norm{A(z)}_\F^2\le r\norm z^2.
\end{equation}
The second assertion follows by summing $\norm{B_a}_\F^2=1$; the third follows from
$\sum_a\norm{B_a^{\mathsf T}z}^2\le\sum_a\norm{B_a}_\F^2\norm z^2$.
Linearity of $A$ makes $z\mapsto\norm{A(z)}$ $1$-Lipschitz and
$z\mapsto\norm{A(z)}_\F$ $\sqrt r$-Lipschitz.

\begin{lemma}[Projected-row moments]\label{lem:row}
For every real $p\ge2$,
\begin{equation}\label{eq:row}
 \boxed{\bigl\|\norm y\bigr\|_{L_p}^{\,2}
 \le r(2p-3)+(p-2)^2.}
\end{equation}
In particular $\|\norm y\|_{L_p}\le\sqrt{2rp}+p$.
\end{lemma}
\begin{proof}
\step{1. Bound the moments of the two Lipschitz norms of $A(g)$.}
By \eqref{eq:collective}, $z\mapsto\norm{A(z)}_\F$ is $\sqrt r$-Lipschitz with
squared $L_2$ norm $r$, and $z\mapsto\norm{A(z)}$ is $1$-Lipschitz with squared
$L_2$ norm at most $r$. Lemma~\ref{lem:lip} gives
\[
 \bigl\|\norm{A(g)}_\F\bigr\|_{L_p}^2\le r(p-1),\qquad
 \bigl\|\norm{A(g)}\bigr\|_{L_p}^2\le r+p-2.
\]

\step{2. Condition on $g$ and integrate over $h$.}
For fixed $g$, the function $h\mapsto\norm{A(g)h}$ has Lipschitz constant
$\norm{A(g)}$ and squared $L_2$ norm $\norm{A(g)}_\F^2$. Lemma~\ref{lem:lip},
applied in the variable $h$ alone, gives
\[
 (\E_h\norm{A(g)h}^p)^{2/p}
 \le\norm{A(g)}_\F^2+(p-2)\norm{A(g)}^2.
\]

\step{3. Combine the two factor bounds.}
Take $L_{p/2}$ norms in $g$ on both sides. Minkowski's inequality in $L_{p/2}$
separates the two summands on the right, and Step~1 bounds each of them:
\[
 \bigl\|\norm y\bigr\|_{L_p}^2
 \le r(p-1)+(p-2)(r+p-2)
 =r(2p-3)+(p-2)^2.
\]
All coefficients remain collected. No common Schmidt basis and no independence of
the coordinates of $y$ is assumed.
\end{proof}

\paragraph{Relation to Banach-space Gaussian chaos bounds.}
The order $\|\norm y\|_{L_p}\lesssim\sqrt{rp}+p$ is already a consequence of
Adamczak--Lata\l a--Meller~\cite[Corollary~12]{ALM}, with their function-space
exponent set to $2$. To check the specialization, set
$c_{ij}=U^{\mathsf T}(e_i\otimes e_j)$ and encode the bilinear form using the
symmetric coefficient array with cross blocks $c_{ij}/2$ and zero diagonal
blocks, acting on the concatenated Gaussian $(g,h)$. The four coefficient
quantities in that corollary are bounded, in their displayed order, by
$\sqrt{r/2}$, $\sqrt r/2$, $1/\sqrt2$, and $1/2$.
This gives the stated order. Lemma~\ref{lem:row} supplies the explicit bound
\eqref{eq:row}, not a claim that this moment order is new. The cited corollary
is an antecedent for the shared local estimate, not by itself the centered
sample-covariance theorem.

\subsection{Directional fourth moments and matrix variance}
The normal fourth-moment calculation below is a special case of Isserlis's
product-moment formula~\cite{Isserlis}, written out to fix the constant $9$.
For the related rank-one probing model, see also~\cite{BK}.
\begin{lemma}[Fourth moments]\label{lem:four}
For every $x\in\R^r$,
\begin{equation}\label{eq:four}
 \E\ip{x}{y}^4\le9\norm x^4.
\end{equation}
Consequently, with $Y=yy^{\mathsf T}$ and $X=Y-I_r$,
\begin{equation}\label{eq:variance}
 \E Y^2\preceq9rI_r,\qquad
 \E X^2\preceq v_r I_r,\qquad v_r:=9r-1.
\end{equation}
\end{lemma}
\begin{proof}
\step{1. The directional fourth moment.}
For unit $x$, put $B_x=\sum_a x_aB_a$. Equation~\eqref{eq:B-orthonormal} gives
$\norm{B_x}_\F=1$. Conditional on $g$, $g^{\mathsf T}B_xh$ is Gaussian with
variance $g^{\mathsf T}Qg$, where $Q=B_xB_x^{\mathsf T}\succeq0$ and $\tr Q=1$.
Diagonalizing $Q$ gives
\[
 \E\ip{x}{y}^4=3\E(g^{\mathsf T}Qg)^2
 =3(\tr Q)^2+6\tr Q^2\le9.
\]
Homogeneity proves \eqref{eq:four}.

\step{2. The matrix second moment.}
For unit $x$, Cauchy--Schwarz in each summand gives
\[
 x^{\mathsf T}(\E Y^2)x
 =\sum_{a=1}^r\E[y_a^2\ip{x}{y}^2]
 \le\sum_{a=1}^r(\E y_a^4)^{1/2}(\E\ip{x}{y}^4)^{1/2}
 \le9r.
\]

\step{3. Center.}
Finally $\E X^2=\E Y^2-2\E Y+I_r=\E Y^2-I_r$.
\end{proof}
The common-factor case $y=\eta\xi$ has
$\|\norm y\|_{L_p}=\|\eta\|_{L_p}\|\norm\xi\|_{L_p}$, of order
$\sqrt{rp}+p$. It calibrates the uniform order of Lemma~\ref{lem:row}, but is not
used to restrict its upper-bound hypotheses.

\clearpage
\section{Chain A: matrix moments, truncation, and concentration}\label{sec:A}
\begin{takeaway}
\textbf{Baseline A.} The constants in this chain remain $9r$ and $72(r+s)$.
The Wick-assisted concentration variant is separated in Appendix~\ref{app:A-wick}.
Chain A never constructs a graded multiplication operator. It converts the common
row estimates into operator-valued moments of one rank-one matrix, truncates only
for the exponential series, and uses matrix Laplace concentration. The final
statement is for the original, unmodified sketch.
\end{takeaway}

The overall architecture has direct precedents in the Khatri--Rao literature
identified in Section~\ref{sec:attribution-intro}. Reproving the concentration
step below does not make that step, or the use of a truncated sample covariance,
a new general method. The displayed constants belong to the implementation in
this note.

\subsection{A1. Keep matrix orientation when estimating powers}
\begin{lemma}[Explicit one-row matrix moments]\label{lem:raw-moments}
For $Y=yy^{\mathsf T}$ and every integer $p\ge3$,
\begin{align}
 \E Y^p
 &\preceq 3\bigl[(8p-11)r+(4p-6)^2\bigr]^{p-1}I_r,\label{eq:raw-moment}\\
 \E Y^p
 &\preceq\frac{p!}{2}\,9r\,[72(r+p)]^{p-2}I_r.\label{eq:72-moment}
\end{align}
For $p=2$, $\E Y^2\preceq9rI_r$.
\end{lemma}
\begin{proof}
\step{1. Bound each directional moment of the matrix power.}
The rank-one identity
\[
 Y^p=\norm y^{2p-2}yy^{\mathsf T}
\]
lets us keep the direction $x$ until after expectation. For any unit $x$,
\begin{align*}
 x^{\mathsf T}(\E Y^p)x
 &=\E\bigl[\norm y^{2p-2}\ip{x}{y}^2\bigr]\\
 &\le(\E\norm y^{4p-4})^{1/2}(\E\ip{x}{y}^4)^{1/2}\\
 &\le3\bigl\|\norm y\bigr\|_{L_{4p-4}}^{2p-2}\\
 &\le3\bigl[(8p-11)r+(4p-6)^2\bigr]^{p-1}.
\end{align*}
The second line is Cauchy--Schwarz; the third bounds the directional factor by
Lemma~\ref{lem:four}; the last applies Lemma~\ref{lem:row} at exponent $4p-4$.
Since $x$ is arbitrary, this proves \eqref{eq:raw-moment} as a Loewner bound.

\step{2. Pass to the uniform moment scale $72(r+p)$.}
The scalar inequality
\begin{equation}\label{eq:72-scalar}
 3\bigl[(8p-11)r+(4p-6)^2\bigr]^{p-1}
 \le\frac{p!}{2}\,9r\,[72(r+p)]^{p-2}
\end{equation}
is proved for every $r\ge1,p\ge3$ in Appendix~\ref{app:72}, giving
\eqref{eq:72-moment}. The $p=2$ case is Lemma~\ref{lem:four}.
\end{proof}
The distinction is between $\norm{\E Y^p}$ and $\E\norm{Y}^p$.
Replacing $Y$ by its norm before taking expectation discards the direction used above.

\subsection{A2. Truncate without paying the full cutoff height}
The untruncated $Y$ need not have a finite positive exponential moment. For example,
when $r=1$ and $y=gh$, $Y=g^2h^2$ has infinite moment-generating function at every
positive argument. Conditional on $g$, the Gaussian integral in $h$ already diverges
whenever $2\theta g^2\ge1$, an event of positive probability.
Thus a positive-matrix exponential theorem cannot be applied to $Y$ directly.

\begin{lemma}[All-order moments of a bounded auxiliary matrix]\label{lem:truncate}
Fix an integer $s\ge2$ and define
\begin{equation}\label{eq:truncation}
 R_s=72(r+s),\qquad L_s=sR_s,\qquad
 W=Y\one_{\{\norm y^2\le L_s\}},\qquad\Sigma=\E W.
\end{equation}
Then
\begin{align}
 \Prob\{W\ne Y\}&\le e^{-2s},\label{eq:clip}\\
 0\preceq I_r-\Sigma&\preceq3e^{-s}I_r,\label{eq:bias}\\
 \E W^p&\preceq\frac{p!}{2}\,9r\,R_s^{p-2}I_r
 \qquad\text{for every integer }p\ge2.\label{eq:all-order}
\end{align}
\end{lemma}
\begin{proof}
\step{1. Bound the event on which the sample is changed.}
At exponent $2s$, Lemma~\ref{lem:row} gives
\[
 \bigl\|\norm y^2\bigr\|_{L_s}
 \le r(4s-3)+(2s-2)^2\le4s(r+s).
\]
Markov's inequality therefore gives
\[
 \Prob\{\norm y^2>L_s\}
 \le\left(\frac{4s(r+s)}{72s(r+s)}\right)^s
 =18^{-s}\le e^{-2s}.
\]

\step{2. Charge the covariance bias in each direction.}
For a unit vector $x$, Lemma~\ref{lem:four} and Cauchy--Schwarz imply
\[
 x^{\mathsf T}(I_r-\Sigma)x
 =\E\bigl[\ip{x}{y}^2\one_{\{\norm y^2>L_s\}}\bigr]
 \le3\Prob\{\norm y^2>L_s\}^{1/2}\le3e^{-s}.
\]
The integrand is nonnegative, proving both sides of \eqref{eq:bias}.

\step{3. Prove the matrix moment bound through order $s$.}
For $p=2$, $\E W^2\preceq\E Y^2\preceq9rI_r$.
For $3\le p\le s$, the scalar event in \eqref{eq:truncation} gives
$W^p\preceq Y^p$ pointwise. Lemma~\ref{lem:raw-moments} and $p\le s$ yield
\[
 \E W^p\preceq\frac{p!}{2}\,9r\,[72(r+p)]^{p-2}I_r
 \preceq\frac{p!}{2}\,9r\,R_s^{p-2}I_r.
\]

\step{4. Extend to every order required by the exponential series.}
For an integer $p>s$, spectral calculus on $0\preceq W\preceq L_sI_r$ gives
\[
 W^p\preceq L_s^{p-s}W^s.
\]
Consequently,
\[
 \E W^p\preceq\frac{s!}{2}\,9r\,R_s^{s-2}L_s^{p-s}I_r
 =\frac{s!s^{p-s}}{2}\,9r\,R_s^{p-2}I_r
 \preceq\frac{p!}{2}\,9r\,R_s^{p-2}I_r.
\]
The last inequality follows by comparing each factor of
$(s+1)(s+2)\cdots p$ with $s$. This proves all of \eqref{eq:all-order},
not just the first $s$ moments.
\end{proof}
\begin{takeaway}
The almost-sure cutoff height is $L_s=sR_s$, but the matrix moment-growth scale is
only $R_s$. Charging $L_s$ as the bounded-summand scale would lose the factor $s$.
\end{takeaway}

\subsection{A3. A matrix Laplace step with valid centering}\label{sec:A-laplace}
The classical matrix-analysis input is Lieb's trace-concavity theorem: for a fixed
self-adjoint $H$, the map
\begin{equation}\label{eq:lieb}
 Q\longmapsto\tr\exp(H+\log Q)
\end{equation}
is concave on the positive-definite cone~\cite{Lieb,Tropp}.
We use no assumption that a random matrix commutes with its expectation.

\begin{lemma}[Matrix Laplace iteration with an offset]\label{lem:laplace}
For independent self-adjoint random matrices $A_1,\ldots,A_m$ with finite matrix
exponential moments, and deterministic self-adjoint $H$,
\begin{equation}\label{eq:laplace}
 \E\tr\exp\left(H+\sum_{i=1}^mA_i\right)
 \le\tr\exp\left(H+\sum_{i=1}^m\log\E e^{A_i}\right).
\end{equation}
\end{lemma}
\begin{proof}
For one random matrix $A$, write $Q=e^A$. Jensen's inequality applied to
\eqref{eq:lieb} gives
\[
 \E\tr\exp(H+A)\le\tr\exp(H+\log\E e^A).
\]
For the sum, condition first on $A_1,\ldots,A_{m-1}$, treating all their terms as
part of $H$. Independence makes $\log\E e^{A_m}$ deterministic. Repeat for
$A_{m-1},\ldots,A_1$. This is the matrix Laplace iteration used by Tropp;
no reordering of products or commuting-mean identity is involved~\cite{Tropp}.
\end{proof}

The next lemma is the positive-raw-moment version of the usual subexponential
matrix Bernstein argument~\cite[Theorem~6.2 and Lemma~6.8]{Tropp}.
Its proof uses the deterministic offset in Lemma~\ref{lem:laplace} so that no
commutation of a sample and its mean is assumed; it is not presented as a new
abstract concentration inequality.

\begin{lemma}[Concentration from positive matrix moments]\label{lem:bernstein}
Let $W_1,\ldots,W_m$ be independent, identically distributed bounded positive
semidefinite $r\times r$ matrices with $\E W_i=\Sigma$. Suppose that for $p\ge2$,
\[
 \E W_i^p\preceq\frac{p!}{2}\,vR^{p-2}I_r.
\]
Then for every $u>0$,
\begin{equation}\label{eq:bernstein}
 \Prob\left\{\left\|\frac1m\sum_iW_i-\Sigma\right\|
 >\sqrt{\frac{2vu}{m}}+\frac{2Ru}{m}\right\}\le2re^{-u}.
\end{equation}
\end{lemma}
\begin{proof}
\step{1. Bound the logarithm of the uncentered exponential moment.}
For $|\theta|R<1$, positivity of each $W_i^p$ and the exponential series imply
\[
 \E e^{\theta W_i}-I_r-\theta\Sigma
 \preceq\sum_{p\ge2}\frac{|\theta|^p}{p!}\E W_i^p
 \preceq\frac{\theta^2v}{2(1-|\theta|R)}I_r.
\]
For positive definite $Q$, spectral calculus gives $\log Q\preceq Q-I_r$.
Therefore
\begin{equation}\label{eq:log-mgf}
 \log\E e^{\theta W_i}
 \preceq\theta\Sigma+
 \frac{\theta^2v}{2(1-|\theta|R)}I_r.
\end{equation}

\step{2. Center by a deterministic offset inside the trace exponential.}
In Lemma~\ref{lem:laplace}, take $A_i=\theta W_i$ and $H=-\theta m\Sigma$.
Using \eqref{eq:log-mgf},
\[
 \E\tr\exp\left(\theta\sum_i(W_i-\Sigma)\right)
 \le r\exp\left(\frac{m\theta^2v}{2(1-|\theta|R)}\right).
\]
Here a Loewner bound on the exponent gives a bound on each eigenvalue and hence
on the trace exponential. No operator monotonicity of the exponential is needed.

\step{3. Convert the exponential bound to the two spectral tails.}
For a self-adjoint $D$ and $\theta>0$, the event $\lambda_{\max}(D)>x$ implies
$\tr e^{\theta D}>e^{\theta x}$. Markov, followed by
$\theta=x/(mv+Rx)$, gives
\[
 \Prob\left\{\lambda_{\max}\left(\sum_i(W_i-\Sigma)\right)>x\right\}
 \le r\exp\left(-\frac{x^2}{2(mv+Rx)}\right).
\]
Use negative $\theta$ for the lower spectral tail and add the probabilities.
Finally $x=\sqrt{2mvu}+2Ru$ satisfies $x^2\ge2u(mv+Rx)$, proving
\eqref{eq:bernstein}.
\end{proof}

\subsection{A4. Restore the original samples and solve for \texorpdfstring{$m$}{m}}
\begin{theorem}[Explicit Chain A certificate]\label{thm:A}
Let $t=\log(4r/\delta)$ and $s=\lceil\log(4m/\delta)\rceil$. Then
\begin{equation}\label{eq:A-tail}
 \boxed{\Prob\left\{\norm Z>
 \sqrt{\frac{18rt}{m}}+\frac{144(r+s)t}{m}+3e^{-s}\right\}\le\delta.}
\end{equation}
In particular, having the displayed allowance at most $\varepsilon$ is an explicit
sufficient condition for the embedding guarantee.
\end{theorem}
\begin{proof}
The integer $s$ is at least $2$.

\step{1. Replace every sample by its truncation.}
Apply Lemma~\ref{lem:truncate} separately to the
independent samples, with the same deterministic cutoff. A union bound gives
\[
 \Prob\{\exists i:W_i\ne Y_i\}\le me^{-2s}
 \le\frac{\delta^2}{16m}\le\frac\delta{16}.
\]

\step{2. Concentrate the truncated average.}
Lemma~\ref{lem:bernstein}, with $v=9r$, $R=R_s=72(r+s)$, and $u=t$, fails with
probability at most $2re^{-t}=\delta/2$.

\step{3. Return to the original sketch.}
On the intersection of the two good events, every $W_i$ equals $Y_i$, so
\begin{align*}
 \norm Z
 &\le\left\|\frac1m\sum_iW_i-\Sigma\right\|+\norm{\Sigma-I_r}\\
 &\le\sqrt{\frac{18rt}{m}}+\frac{144(r+s)t}{m}+3e^{-s}.
\end{align*}
The total failure probability is less than $\delta$. The conclusion concerns the
original $S$, not an estimator that clips its rows during execution.
\end{proof}

\begin{proposition}[The common sample order follows from Chain A]\label{prop:A-size}
The explicit certificate \eqref{eq:A-tail} implies \eqref{eq:main-size} with a
universal constant.
\end{proposition}
\begin{proof}
\step{1. Simplify the allowance to two monotone terms.}
Since $s\le\log(4m/\delta)+1$ and $3e^{-s}\le3\delta/(4m)$, the allowance in
\eqref{eq:A-tail} is at most
\begin{equation}\label{eq:A-rough}
 C_1\left[\sqrt{\frac{rt}{m}}+
 \frac{(r+\log(4m/\delta)+1)t}{m}\right]
\end{equation}
for an absolute $C_1$. Both terms decrease with real $m\ge1$; for the second,
differentiate $(r+\log(4m/\delta)+1)/m$.

\step{2. Evaluate at the candidate sample size.}
Set
\[
 m_0=K\left(\frac{rt}{\varepsilon^2}+\frac{t^2}{\varepsilon}\right)
 =K\varepsilon^{-2}t(r+\varepsilon t),\qquad K\ge1.
\]
Here $t>1$, $r\le e^t$, and $t\le e^t$, so
\[
 \log(4m_0/\delta)+1
 \le\log K+2\log(1/\varepsilon)+3t+\log2+1.
\]

\step{3. Bound each term and fix the constant.}
First,
\[
 \sqrt{\frac{rt}{m_0}}\le\frac\varepsilon{\sqrt K}.
\]
Second,
\begin{align*}
 \frac{(r+\log(4m_0/\delta)+1)t}{m_0}
 &\le\frac{\varepsilon^2}{K}
 \frac{r+3t+2\log(1/\varepsilon)+\log K+\log2+1}{r+\varepsilon t}\\
 &\le\frac{C_2(1+\log K)}K\,\varepsilon.
\end{align*}
For the last step, use $r\ge1$, $r+\varepsilon t\ge\varepsilon t$, and
$\varepsilon\log(1/\varepsilon)\le1/e$.
Choose $K$ so that $C_1[K^{-1/2}+C_2(1+\log K)/K]\le1$.
Monotonicity of \eqref{eq:A-rough} handles every larger integer $m$.
\end{proof}

\clearpage
\section{Chain B: exact graded multiplication and vacuum propagation}\label{sec:B}
\begin{takeaway}
\textbf{Baseline B.} Retain $v_r=9r-1$ and the original $B_q(r;\lambda)$.
The Wick--Hermite refinements belong to the separate Chain C.
Chain B uses the four core stages of the framework: exact multiplication,
reachable graded states, coefficient-sensitive contractions, and vacuum-moment
propagation. No row is clipped and no matrix Laplace theorem is invoked.
The central estimate charges each already-active sample its own local degree
\emph{before} summing over samples.
\end{takeaway}

\subsection{B1. The original-law multiplication identity}
Multiplication in $L_2$ and orthogonal-polynomial coordinates are classical;
Wiener~\cite{Wiener} is a source for homogeneous Gaussian chaos, and
Accardi--Bo\.{z}ejko~\cite{AccardiBozejko} provide a probability-law/creation-operator
representation precedent. We use these as background attributions, not as claims
that either source proves the present embedding theorem. The identity below is
simply Parseval applied to multiplication.
Let
\[
 \mathcal H=\R^r\otimes\bigotimes_{i=1}^m
 L_2(\gamma_{n_1}\otimes\gamma_{n_2}),\qquad \Omega=1.
\]
The external factor $\R^r$ carries the matrix indices, one $L_2$ factor carries
the randomness of one sample, and the vacuum $\Omega$ is the constant function
$1$ in every sample register.
Define $\mathcal M_Zf=Zf$ on the vector polynomials in this Hilbert space.
For an external standard basis vector $e_a$,
\[
 \mathcal M_Z^j(e_a\otimes\Omega)(g_1,h_1,\ldots,g_m,h_m)=Z^je_a.
\]
All of these polynomials are square integrable. The pointwise matrix $Z$ is
symmetric because each $Y_i=y_i y_i^{\mathsf T}$ is symmetric. Thus
$\mathcal M_Z$ is symmetric on its invariant polynomial domain; we do not
assert self-adjointness of this unbounded restriction. The identity follows
directly from $e_a^{\mathsf T}Z^{2q}e_a=\norm{Z^qe_a}^2$ pointwise:
\[
 \E\,e_a^{\mathsf T}Z^{2q}e_a
 =\ip{e_a\otimes\Omega}{\mathcal M_Z^{2q}(e_a\otimes\Omega)}_{\mathcal H}
 =\norm{\mathcal M_Z^q(e_a\otimes\Omega)}_{\mathcal H}^2.
\]
Summing over $a$ gives the identity the chain rests on:
\begin{equation}\label{eq:vacuum}
 \boxed{\E\tr Z^{2q}
 =\sum_{a=1}^r\norm{\mathcal M_Z^q(e_a\otimes\Omega)}_{\mathcal H}^2.}
\end{equation}
The factor $r$ that will eventually arise counts external vacuum seeds. We do not
take a trace over the generally much larger Hermite occupation space.
Every visit to sample $i$ still uses the same $g_i,h_i$.

\subsection{B2. Hermite degree, active samples, and the exact cutoff}
Let $h_\ell=\operatorname{He}_\ell/\sqrt{\ell!}$ denote the normalized probabilists'
Hermite polynomial. Its recurrence is
\begin{equation}\label{eq:Hermite}
 xh_\ell(x)=\sqrt{\ell+1}\,h_{\ell+1}(x)+\sqrt\ell\,h_{\ell-1}(x),
 \qquad h_{-1}=0.
\end{equation}
Products of these functions form an orthonormal basis in each Gaussian sample
space. Let $N_i$ record the total local degree in $(g_i,h_i)$.
Let $P_i$ be the orthogonal projection onto the local constant function, tensored
with identities on all other factors, and put $Q_i=I-P_i$.

\begin{definition}[An active sample]
A sample is active in a product Hermite component when its local factor is
nonconstant, equivalently when it lies in $Q_iL_2$. The exact active set is the
set of such indices. Its size is a second grading, distinct from total degree.
\end{definition}
An active set is not a history of labels that have ever occurred. A sample can
return to its local vacuum, and then be activated again by a later multiplication.

There are exact sign-flip symmetries. Replacing the entire vector $g_i$ by $-g_i$
changes $y_i$ to $-y_i$ but leaves $Y_i$, $X_i$, and $Z$ unchanged; the same holds
for $h_i$. We restrict to functions having even total Hermite degree separately
in every $g_i$ and every $h_i$ register. The vacuum belongs to this subspace,
and multiplication by $Z$ preserves it.

After $j$ multiplications from vacuum,
\begin{equation}\label{eq:reachability}
 \sum_iN_i\le4j,\qquad \#\{\text{active samples}\}\le j.
\end{equation}
Indeed, $X_i$ has polynomial degree at most four and acts on only one sample.

Let $\Pi_q$ project onto the symmetry subspace with total degree at most $4q$
and at most $q$ active samples, and set
\begin{equation}\label{eq:cutoff}
 M_q=\Pi_q\mathcal M_Z\Pi_q.
\end{equation}
This finite-dimensional compression is symmetric, hence self-adjoint.
Inducting on $j$ using
\eqref{eq:reachability} gives
\begin{equation}\label{eq:exact-cutoff}
 M_q^j(e_a\otimes\Omega)=\mathcal M_Z^j(e_a\otimes\Omega),\qquad 0\le j\le q.
\end{equation}
Thus \eqref{eq:vacuum} remains exact with $M_q$ in place of $\mathcal M_Z$.
This is the elementary finite-moment compression principle as organized in the
framework~\cite{Framework}; its proof here is the reachability induction just
given. Neither that induction nor the underlying polynomial representation is
claimed as a new general theorem. No bound on the full unbounded Gaussian
multiplication operator is required.

\subsection{B3. Split activation, return, and repeated visits}
Write $\mathcal X_i$ for multiplication by $X_i$. Centering is the exact identity
$P_i\mathcal X_iP_i=0$. Therefore
\[
 \mathcal X_i=Q_i\mathcal X_iP_i+P_i\mathcal X_iQ_i+Q_i\mathcal X_iQ_i.
\]
Define, on the common cutoff,
\begin{equation}\label{eq:C-R}
 C=\Pi_q\sum_iQ_i\mathcal X_iP_i\Pi_q,\qquad
 \mathcal R=\Pi_q\sum_iQ_i\mathcal X_iQ_i\Pi_q.
\end{equation}
Then
\begin{equation}\label{eq:decomposition}
 \boxed{M_q=\frac1m(C+C^*+\mathcal R).}
\end{equation}
The operator $C$ increases active-sample count by one; $C^*$ decreases it by one;
$\mathcal R$ preserves the exact active set. These are actual full-space
compressions, so using adjointness is legitimate.
Let $\mathcal K_k$ denote the cutoff sector with exactly $k$ active samples.

\subsubsection{Activation and return are controlled by matrix variance}
\begin{lemma}[Fresh-sample contraction]\label{lem:fresh}
With $v_r=9r-1$, for $0\le k<\min(q,m)$,
\begin{equation}\label{eq:fresh}
 \norm{C:\mathcal K_k\to\mathcal K_{k+1}}
 \le\sqrt{(k+1)(m-k)v_r}.
\end{equation}
The reverse map $C^*$ has the same bound.
\end{lemma}
\begin{proof}
\step{1. Compute the single-insertion Gram before the output cutoff.}
For a sample currently in its local vacuum,
\begin{align*}
 (Q_i\mathcal X_iP_i)^*(Q_i\mathcal X_iP_i)
 &=P_i\mathcal X_iQ_i\mathcal X_iP_i\\
 &=P_i\mathcal X_i^2P_i
 = (\E X_i^2)P_i\preceq v_rP_i.
\end{align*}
The middle equality uses $P_i\mathcal X_iP_i=0$.
It remains valid with arbitrary Hilbert-valued coefficients in the other sample
registers. Output projection by $\Pi_q$ can only decrease a norm.

\step{2. Count the inputs that can meet at one output set.}
Decompose $f=\sum_{|S|=k}f_S$ by exact active set; these inputs are orthogonal.
For an output set $T$ of size $k+1$, the contributing inputs are
$f_{T\setminus\{i\}}$, one for each $i\in T$. Cauchy--Schwarz gives
\[
 \left\|\sum_{i\in T}\Pi_q Q_i\mathcal X_iP_i f_{T\setminus\{i\}}\right\|^2
 \le(k+1)v_r\sum_{i\in T}\norm{f_{T\setminus\{i\}}}^2.
\]
Sum over $T$. Each input set $S$ appears once for each of the $m-k$ absent
labels, yielding
\[
 \norm{Cf}^2\le(k+1)(m-k)v_r\sum_{|S|=k}\norm{f_S}^2.
\]
Take square roots and then use adjointness.
\end{proof}

\subsubsection{Repeated visits are controlled by local degree}
For one sample define the row multiplier
\[
 Tf=y^{\mathsf T}f,\qquad T^*T=\mathcal M_Y
\]
as a quadratic-form identity on polynomials. All following estimates allow a
Hilbert-valued coefficient space for the other registers. H\"older is applied
only in the current sample variables, with that coefficient space kept intact.
Before the finite cutoff, identities involving multipliers and their adjoints
are understood as polynomial quadratic-form identities; after compression the
adjoints are ordinary finite-dimensional ones.

\begin{lemma}[Local degree-dependent row contraction]\label{lem:local}
Let $f_d$ have homogeneous local Hermite degree $d\ge1$. For every real
$\lambda>0$,
\begin{equation}\label{eq:local}
 \norm{Tf_d}_2^2\le
 e^{2/\lambda}\bigl[(4\lambda d+1)r+4\lambda^2d^2\bigr]\norm{f_d}_2^2.
\end{equation}
For a finite sum $f=\sum_{d\ge1}f_d$ with no local vacuum component,
\begin{equation}\label{eq:local-sum}
 \norm{Tf}_2^2\le3e^{2/\lambda}\sum_{d\ge1}
 \bigl[(4\lambda d+1)r+4\lambda^2d^2\bigr]\norm{f_d}_2^2.
\end{equation}
\end{lemma}
\begin{proof}
\step{1. Choose the H\"older exponent from the actual local degree.}
Set $p=\lambda d+1>1$. H\"older and \eqref{eq:HC} give
\begin{align*}
 \norm{y^{\mathsf T}f_d}_2
 &\le\bigl\|\norm y\bigr\|_{L_{2p}}\,\norm{f_d}_{L_{2p/(p-1)}}\\
 &\le\bigl\|\norm y\bigr\|_{L_{2p}}
 \left(1+\frac2{\lambda d}\right)^{d/2}\norm{f_d}_2.
\end{align*}
Lemma~\ref{lem:row} bounds the square of the first factor by
\[
 r(4p-3)+(2p-2)^2=(4\lambda d+1)r+4\lambda^2d^2.
\]
Since $(1+2/(\lambda d))^d\le e^{2/\lambda}$, this proves
\eqref{eq:local}. The multiplier and the polynomial use the same sample;
no independence between $y$ and $f_d$ is assumed.

\step{2. Sum homogeneous inputs using their output-degree overlap.}
Every coordinate of $y$ is a pure second-chaos polynomial. In fact it is a linear
combination of $g_a h_b$, with distinct factor coordinates. Write $\mathcal F_e$
for the closed span of the components of exact local Hermite degree $e$, with
arbitrary coefficients in the other registers; set $\mathcal F_e=\{0\}$ for
$e<0$. Applying the Hermite recurrence \eqref{eq:Hermite} twice shows
\[
 Tf_d\in\mathcal F_{d-2}\oplus\mathcal F_d\oplus\mathcal F_{d+2}.
\]
At any fixed output degree, at most three input degrees contribute.
Cauchy--Schwarz on those contributions, followed by orthogonality of the output
degrees, gives $\norm{Tf}_2^2\le3\sum_d\norm{Tf_d}_2^2$.
Now use \eqref{eq:local}.
\end{proof}

Define the deterministic allowance
\begin{equation}\label{eq:Bq}
 \boxed{B_q(r;\lambda)=3e^{2/\lambda}
 \bigl[(16\lambda+1)rq+64\lambda^2q^2\bigr].}
\end{equation}
The positive real $\lambda$ is a proof parameter, not a property of the subspace.
In particular,
\begin{equation}\label{eq:B-special}
 B_q(r;2)=99e\,rq+768e\,q^2,\qquad
 B_q(r;1)=51e^2rq+192e^2q^2.
\end{equation}

\begin{lemma}[Sum of already-active sample interactions]\label{lem:repeated}
On the entire cutoff of \eqref{eq:cutoff},
\begin{equation}\label{eq:repeated}
 \norm{\mathcal R}\le B_q(r;\lambda).
\end{equation}
Its restriction to the vacuum sector is zero.
\end{lemma}
\begin{proof}
\step{1. Fix the exact active set and expose the positive part.}
On a support sector $S$ with $k=|S|\le q$, the quadratic form is
\begin{equation}\label{eq:R-form}
 \ip f{\mathcal Rf}=\sum_{i\in S}\norm{T_if}_2^2-k\norm f_2^2.
\end{equation}
This follows because $Q_if=f$ for $i\in S$ and $Q_if=0$ otherwise, together
with the one-sample identity
$\ip f{\mathcal X_if}=\ip f{(\mathcal M_{Y_i}-I)f}=\norm{T_if}_2^2-\norm f_2^2$.
Thus $\mathcal R\succeq-kI$ on this sector.

\step{2. Sum each label's actual degree before relaxing any budget.}
Let $f=\sum_{\mathbf d}f_{\mathbf d}$ be the joint local-degree decomposition.
The number operators $N_i$ commute, and their joint eigenspaces are orthogonal.
Apply \eqref{eq:local-sum} to each $i\in S$ and then combine the sums:
\begin{align*}
 \sum_{i\in S}\norm{T_if}_2^2
 &\le3e^{2/\lambda}\sum_{\mathbf d}
 \left[r\left(4\lambda\sum_{i\in S}d_i+k\right)
       +4\lambda^2\sum_{i\in S}d_i^2\right]\norm{f_{\mathbf d}}_2^2\\
 &\le3e^{2/\lambda}
 \left[(16\lambda+1)rq+64\lambda^2q^2\right]\norm f_2^2.
\end{align*}
The last inequality uses precisely
\[
 \sum_i d_i\le4q,\qquad k\le q,\qquad
 \sum_i d_i^2\le\left(\sum_i d_i\right)^2\le16q^2.
\]
Charging every active label degree $4q$ separately would be an unnecessary extra
factor of $q$; the displayed summation avoids that loss.

\step{3. Control both sides of the centered form.}
Equation~\eqref{eq:R-form} and the last estimate give
$-kI\preceq\mathcal R\preceq B_q I$.
Since $B_q\ge q\ge k$, its operator norm is at most $B_q$.
Compression by the cutoff preserves the form inequalities. Finally sum over
orthogonal exact-support sectors. On the vacuum sector $S=\varnothing$, the
operator vanishes.
\end{proof}

\subsection{B4. Propagate the vacuum rather than the worst state}
Let $K=\min(q,m)$ and define the nonnegative symmetric tridiagonal matrix $G_q$
on indices $0,\ldots,K$ by
\begin{align}
 (G_q)_{00}&=0,\qquad (G_q)_{kk}=B_q(r;\lambda)/m\quad(k\ge1),\nonumber\\
 (G_q)_{k+1,k}&=(G_q)_{k,k+1}
 =\frac{\sqrt{(k+1)(m-k)v_r}}m\quad(0\le k<K).
 \label{eq:G}
\end{align}
If $n(f)=(\norm{P_{\mathcal K_k}f})_{k=0}^K$ is the vector of sector norms,
Lemmas~\ref{lem:fresh} and~\ref{lem:repeated} give
\begin{equation}\label{eq:sector-step}
 n(M_qf)\le G_qn(f)
\end{equation}
componentwise. Since $G_q$ is entrywise nonnegative, iteration from
$n(e_a\otimes\Omega)=e_0$ yields
\begin{equation}\label{eq:stationary}
 \boxed{\E\tr Z^{2q}\le r\norm{G_q^qe_0}_2^2=r(G_q^{2q})_{00}.}
\end{equation}
This is the nonnegative vacuum-majorant rule used in the framework~\cite{Framework}.
The comparison itself follows from block norms and positivity; the local
activation and repeated-visit estimates supply its application-specific entries.
All occurrences of $e_0$ in scalar majorants refer to the active-count vacuum;
it is distinct from the external basis vectors $e_a$.

\subsubsection{A closed-form consequence}
The scalar Jacobi comparison uses the classical relation between orthogonal
polynomials, finite Jacobi matrices, and Gaussian quadrature; see
Golub--Welsch~\cite{GolubWelsch}. The vacuum weights, rather than the unweighted
trace of that Jacobi matrix, are essential. The required finite-moment identity
is proved explicitly below.
Let
\begin{equation}\label{eq:cq-eta}
 c_q=((2q-1)!!)^{1/(2q)},\qquad
 \eta_q=(r/\delta)^{1/(2q)}.
\end{equation}
\begin{theorem}[Untruncated closed-form moment bound]\label{thm:B-closed}
For every integer $q\ge1$ and every $\lambda>0$,
\begin{equation}\label{eq:B-moment}
 \boxed{(\E\tr Z^{2q})^{1/(2q)}
 \le r^{1/(2q)}\left[c_q\sqrt{\frac{v_r}{m}}+\frac{B_q(r;\lambda)}m\right].}
\end{equation}
Consequently
\begin{equation}\label{eq:B-tail}
 \Prob\left\{\norm Z>
 \eta_q\left[c_q\sqrt{\frac{v_r}{m}}+\frac{B_q(r;\lambda)}m\right]\right\}
 \le\delta.
\end{equation}
\end{theorem}
\begin{proof}
\step{1. Dominate the sector matrix entrywise.}
Let $J_q$ be the $(q+1)\times(q+1)$ Hermite Jacobi matrix with zero diagonal and
$(J_q)_{k+1,k}=(J_q)_{k,k+1}=\sqrt{k+1}$.
After padding $G_q$ by zeros if $m<q$, it is entrywise at most
\[
 bI+aJ_q,\qquad a=\sqrt{v_r/m},\quad b=B_q(r;\lambda)/m,
\]
because $(k+1)(m-k)\le(k+1)m$ and $B_q\ge0$ on the diagonal.
The entries are nonnegative, so matrix powers preserve this comparison, and in
particular $\norm{G_q^qe_0}\le\norm{(bI+aJ_q)^qe_0}$.

\step{2. Identify the comparison vacuum moment as a scalar Gaussian moment.}
By \eqref{eq:Hermite}, $J_q$ is multiplication by a scalar standard Gaussian
$\xi$ compressed to Hermite degrees $0,\ldots,q$. Each application of
$b+a\xi$ raises the degree by at most one, so $q$ applications starting from
$h_0=1$ never leave degrees $0,\ldots,q$, and the compression changes nothing.
Thus
\[
 \norm{(bI+aJ_q)^qe_0}^2=\E(b+a\xi)^{2q}.
\]

\step{3. Separate the two rates and apply Markov.}
Minkowski in scalar $L_{2q}$ gives
\[
 \norm{G_q^qe_0}^{1/q}\le b+a\norm\xi_{L_{2q}}=b+ac_q.
\]
Combine this with \eqref{eq:stationary} to obtain \eqref{eq:B-moment}. Finally,
$\Prob\{\norm Z>x\}\le x^{-2q}\E\tr Z^{2q}$ at the displayed threshold proves
the tail bound \eqref{eq:B-tail}, since $\eta_q^{2q}=r/\delta$.
\end{proof}
The scalar Gaussian $\xi$ describes only the comparison Jacobi operator. It does
not replace any original tensor factor or any repeated sample in the original law.

\begin{corollary}[The common sample order follows from Chain B]\label{cor:B-size}
The framework proof implies \eqref{eq:main-size}. An explicit sufficient row
count for any chosen $q\ge1,\lambda>0$ is
\begin{equation}\label{eq:B-explicit-m}
 m\ge\left\lceil\left(
 \frac{a_q+\sqrt{a_q^2+4\varepsilon b_q}}{2\varepsilon}
 \right)^2\right\rceil,
 \qquad
 a_q=\eta_qc_q\sqrt{v_r},\quad b_q=\eta_qB_q(r;\lambda).
\end{equation}
\end{corollary}
\begin{proof}
\step{1. The exact threshold.}
Equation~\eqref{eq:B-explicit-m} solves
$a_q/\sqrt m+b_q/m\le\varepsilon$ exactly as a quadratic inequality in
$\sqrt m$.

\step{2. The sample-size order.}
Take $q=\lceil\log(r/\delta)\rceil\ge1$ and $\lambda=2$.
Then $\eta_q\le\sqrt e$, $c_q\le\sqrt{2q}$, and $q\le2t$ for
$t=\log(4r/\delta)>1$.
The allowance is therefore bounded by an absolute constant times
\[
 \sqrt{rq/m}+rq/m+q^2/m.
\]
Making each term at most a fixed fraction of $\varepsilon$ requires the sample
order
$\varepsilon^{-2}rt+\varepsilon^{-1}rt+\varepsilon^{-1}t^2$.
Since $\varepsilon<1$, the middle term is absorbed by the first, which gives
\eqref{eq:main-size}.
\end{proof}

\subsection{A stronger certificate: use the budget at each stage}\label{sec:stagewise}
The fixed cutoff at degree $4q$ is exact, but its largest repeated-visit cost need
not be charged at every earlier step. A prefix of length $j-1$ only needs the next
cutoff at degree $4j$ and at most $j$ active samples.

Fix the terminal $q$ and $K=\min(q,m)$. Define a symmetric nonnegative matrix
$G^{[j]}$ on $0,\ldots,K$ for $1\le j\le q$ by
\begin{align}
 (G^{[j]})_{00}&=0,\qquad (G^{[j]})_{kk}=B_j(r;\lambda)/m\quad(k\ge1),\nonumber\\
 (G^{[j]})_{k+1,k}&=(G^{[j]})_{k,k+1}
 =\frac{\sqrt{(k+1)(m-k)v_r}}m.
 \label{eq:stage-matrix}
\end{align}
Entries beyond currently reachable sectors are harmless padding.
Propagate in the stated order:
\begin{equation}\label{eq:ordered-propagation}
 w^{(0)}=e_0,\qquad w^{(j)}=G^{[j]}w^{(j-1)}.
\end{equation}

\begin{theorem}[Stage-dependent vacuum certificate]\label{thm:stagewise}
For the original independent Gaussian law,
\begin{equation}\label{eq:stage-moment}
 \boxed{\E\tr Z^{2q}\le
 r\norm{G^{[q]}G^{[q-1]}\cdots G^{[1]}e_0}_2^2.}
\end{equation}
In particular, the computable inequality
\begin{equation}\label{eq:stage-tail}
 r\norm{G^{[q]}\cdots G^{[1]}e_0}_2^2\le\delta\varepsilon^{2q}
\end{equation}
is sufficient for $\Prob\{\norm Z>\varepsilon\}\le\delta$.
At fixed $q,\lambda,r,m$, this certificate is no worse than the stationary
certificate \eqref{eq:stationary}.
\end{theorem}
\begin{proof}
\step{1. One exact stage.}
For one external seed, put $f_j=\mathcal M_Z^j(e_a\otimes\Omega)$.
By \eqref{eq:reachability}, $f_{j-1}$ belongs to the cutoff with degree
$4(j-1)$ and active count at most $j-1$, while $f_j$ lies in the next cutoff.
Let $M_j=\Pi_j\mathcal M_Z\Pi_j$ use that next cutoff. Then
$f_j=M_jf_{j-1}$ exactly.

\step{2. Propagate sector norms in the stated order.}
Apply Lemmas~\ref{lem:fresh} and~\ref{lem:repeated} with parameter $j$:
\[
 n(f_j)\le G^{[j]}n(f_{j-1}).
\]
Induction gives $n(f_q)\le w^{(q)}$ componentwise. Sum squared norms over all
$r$ external seeds and use \eqref{eq:vacuum} to obtain \eqref{eq:stage-moment}.
Markov then yields \eqref{eq:stage-tail}.

\step{3. Compare with the stationary certificate.}
Finally $B_j(r;\lambda)\le B_q(r;\lambda)$ implies
$G^{[j]}\le G_q$ entrywise. Positivity and induction give
$w^{(q)}\le G_q^qe_0$ componentwise.
\end{proof}

\begin{remark}[Order and energy are essential]
The matrices $G^{[j]}$ generally do not commute. The terminal quantity is the
squared norm of an \emph{ordered product}; it is not a product of top eigenvalues,
a power of an average, or an unordered product of return costs.

To express the argument as an energy, define for a nonnegative sector vector $z$
\[
 \mathcal E_j(z)=\norm{G^{[q]}\cdots G^{[j+1]}z}_2^2,
 \qquad \mathcal E_q(z)=\norm z^2.
\]
The componentwise step inequality and nonnegativity give
$\mathcal E_j(n(f_j))\le\mathcal E_{j-1}(n(f_{j-1}))$.
This is the framework's weighted-energy viewpoint, with a specified finite
certificate rather than a formal re-expression of the unknown exact moment.
\end{remark}

\subsection{An exact limited-independence consequence}
\begin{corollary}[Moment matching of sample blocks]\label{cor:limited}
Fix $q$. Suppose each block $(g_i,h_i)$ has the correct two-factor Gaussian
marginal law, and any set of at most $2q$ distinct blocks has the same joint law
as under full independence. Then $\E\tr Z^{2q}$ is unchanged. All three
framework certificates at this $q$ therefore remain valid.
\end{corollary}
\begin{proof}
Expand $\tr(\sum_iX_i)^{2q}$. Each word uses at most $2q$ distinct block labels.
Its expectation is fixed by their joint law, even if a label is repeated many
times within the word. Sum the finitely many equal expectations and divide by
$m^{2q}$. The operator proof is carried out on the independent reference law;
only its resulting moment is transferred.
\end{proof}
This does not provide a finite-seed Gaussian generator or transfer the lower
bound. No concentration or cutoff argument on the dependent law is asserted.

\clearpage
\section{Chain C: the combined Wick--Hermite framework}\label{sec:C}
\begin{takeaway}[colframe=Amber!80,coltitle=Amber,title={A third proof chain, not a replacement for B}]
Wick contractions determine sharper local radial and covariance bounds.
Those bounds enter B's original-law Hermite multiplication, active-register
cutoff, and ordered vacuum propagation. This hybrid never clips an outcome and
never uses matrix Laplace concentration. Its constants carry a separate
superscript $(C)$; the constants and proofs of A and B remain unchanged.
\end{takeaway}

The starting model is still \eqref{eq:model}. The coefficient matrices $B_a$
are still the arbitrary Frobenius-orthonormal matrices of
\eqref{eq:B-orthonormal}. No common Schmidt basis is assumed in an upper bound.
Wick here means joint Gaussian product-moment contraction~\cite{Isserlis,Janson},
not replacement of ordinary multiplication by the Wick product with its
contraction terms removed. The radialization idea is recorded in
\cite{RMPS}; the particular Khatri--Rao refinement is supplied by
\cite{WickNote}. Both are provenance, not claims of external validation.

\subsection{C1. Contract jointly and retain the sharp radial envelope}\label{c:sec:wick-radial}
For a positive integer $p$, define
\begin{equation}\label{c:eq:wick-parameters}
 \kappa_p=(2p-1)!!,\qquad
 D_{r,p}=\prod_{\ell=0}^{p-1}(r+2\ell),\qquad
 M_p(r)=(\kappa_pD_{r,p})^{1/p}.
\end{equation}
Here $\kappa_p$ counts scalar pairings. It is not the moment root $c_q$ that
appears later in the scalar Jacobi comparison. Set $\kappa_0=D_{r,0}=1$.

\begin{theorem}[Sharp uniform radial moments]\label{c:thm:wick-radial}
For every integer $p\ge1$,
\begin{equation}\label{c:eq:wick-radial}
 \boxed{\E\norm y^{2p}\le \kappa_pD_{r,p}.}
\end{equation}
For every $r$, equality holds on the common-factor subspace $Ue_a=e_a\otimes e_1$. Moreover,
\begin{equation}\label{c:eq:wick-amgm}
 \norm{\norm y}_{L_{2p}}^2\le M_p(r)\le p(r+p-1).
\end{equation}
\end{theorem}
\begin{proof}
\step{1. Contract one scalar projection exactly.}
For $x\in\R^r$, put $B_x=\sum_a x_a B_a$. Then $\norm{B_x}_F^2=\norm x^2$. Conditional Gaussian integration, equivalently Wick's rule, gives
\[
 \E(g^{\mathsf T}B_xh)^{2p}
 =\kappa_p\,\E\left(\sum_j\lambda_j\gamma_j^2\right)^p,
\]
where $\lambda_j\ge0$ are the squared singular values of $B_x$, their sum is $\norm x^2$, and the $\gamma_j$ are independent standard Gaussians. For $x\ne0$, write $w_j=\lambda_j/\norm x^2$, so $\sum_jw_j=1$.
Convexity gives
\[
 \left(\sum_jw_j\gamma_j^2\right)^p
 \le\sum_jw_j\gamma_j^{2p}.
\]
Taking expectations (and handling $x=0$ separately) yields
\begin{equation}\label{c:eq:wick-directional}
 \E\ip{x}{y}^{2p}\le \kappa_p^2\norm x^{2p}.
\end{equation}
For rank-one $B_x$ the inequality is an equality. All repeated Gaussian occurrences were integrated jointly.

\step{2. Recover the radial moment without a coordinatewise norm bound.}
Let $\zeta\sim N(0,I_r)$ be independent of $y$. Conditional on $y$,
\[
 \E_\zeta\ip{\zeta}{y}^{2p}=\kappa_p\norm y^{2p}.
\]
Apply \eqref{c:eq:wick-directional} conditional on $\zeta$, then integrate:
\[
 \kappa_p\E\norm y^{2p}\le \kappa_p^2\E\norm\zeta^{2p}=\kappa_p^2D_{r,p}.
\]
The chi-square moment $D_{r,p}$ follows either from Wick pairing counts or Gaussian polar integration. Divide by $\kappa_p$. This auxiliary-Gaussian argument is equivalent to the spherical averaging in \cite[projected radial moments]{RMPS}; $\zeta$ is a proof device, not an additional sketch factor.

\step{3. Check equality and the elementary envelope.}
For $y=\eta\xi$, with independent $\eta\sim N(0,1)$ and $\xi\sim N(0,I_r)$, both factors in the radial moment separate, giving $\kappa_pD_{r,p}$ exactly. Finally the arithmetic means of the factors in $\kappa_p$ and $D_{r,p}$ are $p$ and $r+p-1$. AM--GM proves \eqref{c:eq:wick-amgm}.
\end{proof}

\begin{remark}[Why this is an improvement in constants, not a new rate]
For $r$ much larger than $p$, the exact envelope is asymptotically $\kappa_p^{1/p}r$, with $\kappa_p^{1/p}\sim2p/e$. The preceding baseline is approximately $4pr$. The local leading ratio therefore approaches $2e$ when both approximations apply. The established order $\sqrt{rp}+p$ has not changed. The exact product also retains the regime in which $p$ is comparable to or larger than $r$.
\end{remark}


\subsection{C2. Keep orthogonal mixed contractions}
A bound on each directional fourth moment alone gives the earlier $9r$. To improve it, keep the orthogonality between distinct subspace coefficient matrices while doing the contractions.

For real matrices $A,B$ of the same shape, define
\begin{equation}\label{c:eq:wick-Q}
 Q(A,B)=\tr(A^{\mathsf T}AB^{\mathsf T}B)
       +\tr(AA^{\mathsf T}BB^{\mathsf T})
       +\tr\big((A^{\mathsf T}B)^2\big).
\end{equation}
The last summand can have either sign. Wick--Isserlis gives the exact identity
\begin{equation}\label{c:eq:mixed-wick}
 \E(g^{\mathsf T}Ah)^2(g^{\mathsf T}Bh)^2
 =\norm A_F^2\norm B_F^2+2\ip{A}{B}_F^2+2Q(A,B).
\end{equation}
One direct verification conditions on $h$. The Gaussian moment in $g$ is
$\norm{Ah}^2\norm{Bh}^2+2(h^{\mathsf T}A^{\mathsf T}Bh)^2$;
expanding its two quadratic-form moments in $h$ gives \eqref{c:eq:mixed-wick}.

\begin{lemma}[Orthogonal mixed fourth moment]\label{c:lem:wick-mixed}
If $\norm A_F=\norm B_F=1$ and $\langle A,B\rangle_F=0$, then
\begin{equation}\label{c:eq:wick-mixed4}
 \boxed{\E(g^{\mathsf T}Ah)^2(g^{\mathsf T}Bh)^2\le4.}
\end{equation}
The constant $4$ in this two-direction statement is attained for suitable
$2\times2$ coefficient matrices.
\end{lemma}
\begin{proof}
\step{1. Use singular coordinates for $A$.}
By orthogonal invariance, $A$ is rectangular diagonal with singular values $s_i\ge0$, $\sum_i s_i^2=1$. Write $b_{ij}$ for the entries of $B$ in these bases. On each off-diagonal pair $(b_{ij},b_{ji})$ inside the diagonal support, the quadratic form $Q$ has matrix
\[
 \begin{pmatrix}s_i^2+s_j^2&s_is_j\\s_is_j&s_i^2+s_j^2\end{pmatrix}.
\]
Its top eigenvalue is $s_i^2+s_j^2+s_is_j\le\tfrac32(s_i^2+s_j^2)\le\tfrac32$. Entries outside such a pair have coefficient at most $1$ (or zero).

\step{2. Use coefficient orthogonality on the diagonal.}
The diagonal contribution is $3\sum_i s_i^2b_{ii}^2$, with $\sum_i s_i b_{ii}=0$. For $u_i=s_i b_{ii}$, its positive and negative parts have equal total mass. Cauchy--Schwarz gives $\sum_i|u_i|\le(\sum_i b_{ii}^2)^{1/2}$. Therefore
\[
 \sum_i u_i^2\le
 \left(\sum_{u_i>0}u_i\right)^2+
 \left(\sum_{u_i<0}|u_i|\right)^2
 \le\frac12\sum_i b_{ii}^2.
\]
The diagonal contribution is at most $\tfrac32\sum_i b_{ii}^2$. Combining all entries gives $Q(A,B)\le\tfrac32\norm B_F^2=\tfrac32$.

\step{3. Finish and check sharpness.}
Equation \eqref{c:eq:mixed-wick} becomes $1+2Q\le4$. Equality holds for
$A=2^{-1/2}\diag(1,1)$ and $B=2^{-1/2}\diag(1,-1)$.
\end{proof}

\begin{corollary}[Uniform matrix second moments]\label{c:cor:wick-variance}
With
\begin{equation}\label{c:eq:variance-parameters}
 \nu_r^{(C)}=4r+5,\qquad v_r^{(C)}=4r+4,
\end{equation}
one has
\begin{equation}\label{c:eq:variance}
 \boxed{\E Y^2\preceq\nu_r^{(C)} I_r,\qquad \E X^2\preceq v_r^{(C)} I_r.}
\end{equation}
\end{corollary}
\begin{proof}
Fix a unit $x$ and rotate the external coefficient basis so that its first vector is $x$. The corresponding first coefficient matrix is $B_x$, and the other $r-1$ coefficient matrices are Frobenius-orthogonal to it. Its own fourth moment is at most $9$, by \eqref{c:eq:wick-directional} with $p=2$; each mixed fourth moment is at most $4$ by Lemma~\ref{c:lem:wick-mixed}. Hence
\[
 x^{\mathsf T}(\E Y^2)x=\E\norm y^2\ip{x}{y}^2\le9+4(r-1)=4r+5.
\]
Finally $\E X^2=\E(Y-I_r)^2=\E Y^2-I_r$.
\end{proof}

\subsubsection{Optional exact covariance information}\label{c:sec:exact-covariance}
The Wick calculation also provides exact directional covariance data. Set
\[
 F_g=\sum_aB_aB_a^{\mathsf T},\qquad F_h=\sum_aB_a^{\mathsf T}B_a.
\]
For unit $x$,
\begin{align}\label{c:eq:wick-exactcov}
 x^{\mathsf T}(\E Y^2)x
 =r+2+2\bigg[&\tr(B_x^{\mathsf T}B_xF_h)+\tr(B_xB_x^{\mathsf T}F_g)\notag\\
 &+\sum_a\tr\big((B_x^{\mathsf T}B_a)^2\big)\bigg].
\end{align}
It follows by summing \eqref{c:eq:mixed-wick}; $\sum_a\ip{B_x}{B_a}_F^2=1$.
An exactly computed or certified upper bound on $\norm{\E X^2}$ can replace $v_r^{(C)}$ in the activation bound below. This does not automatically replace $\nu_r^{(C)}$ in every higher-moment estimate of Chain A.

\begin{remark}[Radial sharpness does not identify the worst covariance]
For the common-factor family $y=\eta\xi$, one has $\E Y^2=3(r+2)I_r$ and $\E X^2=(3r+5)I_r$. That does not justify the universal allowance $3r+5$. For example, let $r=2$ and
\[
 B_1=\diag(a,b),\qquad B_2=\frac1{\sqrt2}\begin{pmatrix}0&1\\1&0\end{pmatrix},
 \qquad a^2+b^2=1,\quad ab=\frac1{12}.
\]
These two matrices are Frobenius-orthonormal. Wick gives
$\E y_1^4=107/12$, $\E y_1^2y_2^2=19/6$, and hence
$e_1^{\mathsf T}(\E X^2)e_1=133/12>11=3r+5$.
The bound $4r+4$ is a valid envelope, not a claimed sharp covariance constant.
\end{remark}


\subsection{C3. Insert the Wick bounds into the Hermite framework}
The probability space, multiplication operator, sign-flip symmetries, cutoff
$\Pi_q$, and compressed matrix $M_q$ are exactly those defined and proved in
Section~\ref{sec:B}. In particular,
\[
 \mathbb E\tr Z^{2q}=\sum_{a=1}^r\|M_q^q(e_a\otimes\Omega)\|^2,
 \qquad mM_q=C+C^*+\mathcal R.
\]
The letter $C$ inside this decomposition denotes the activation operator of
\eqref{eq:C-R}, not the proof-chain label. The operator itself is unchanged;
only its certified norm budgets are sharpened. Repeated visits still refer to
the same original sample. All unbounded multiplier identities below are
polynomial quadratic-form identities; their finite compressions are
self-adjoint finite matrices.

\subsubsection{Activation and return with the improved variance}
\begin{lemma}[Fresh-sample contraction]\label{c:lem:fresh}
With $v_r^{(C)}=4r+4$, for $0\le k<\min(q,m)$,
\begin{equation}\label{c:eq:fresh}
 \norm{C:\mathcal K_k\to\mathcal K_{k+1}}
 \le\sqrt{(k+1)(m-k)v_r^{(C)}}.
\end{equation}
The reverse map $C^*$ has the same bound.
\end{lemma}
\begin{proof}
\step{1. Compute the single-insertion Gram before the output cutoff.}
For a sample currently in its local vacuum,
\begin{align*}
 (Q_i\mathcal X_iP_i)^*(Q_i\mathcal X_iP_i)
 &=P_i\mathcal X_iQ_i\mathcal X_iP_i\\
 &=P_i\mathcal X_i^2P_i
 = (\E X_i^2)P_i\preceq v_r^{(C)}P_i.
\end{align*}
The middle equality uses $P_i\mathcal X_iP_i=0$.
It remains valid with arbitrary Hilbert-valued coefficients in the other sample
registers. Output projection by $\Pi_q$ can only decrease a norm.

\step{2. Count the inputs that can meet at one output set.}
Decompose $f=\sum_{|S|=k}f_S$ by exact active set; these inputs are orthogonal.
For an output set $T$ of size $k+1$, the contributing inputs are
$f_{T\setminus\{i\}}$, one for each $i\in T$. Cauchy--Schwarz gives
\[
 \left\|\sum_{i\in T}\Pi_q Q_i\mathcal X_iP_i f_{T\setminus\{i\}}\right\|^2
 \le(k+1)v_r^{(C)}\sum_{i\in T}\norm{f_{T\setminus\{i\}}}^2.
\]
Sum over $T$. Each input set $S$ appears once for each of the $m-k$ absent
labels, yielding
\[
 \norm{Cf}^2\le(k+1)(m-k)v_r^{(C)}\sum_{|S|=k}\norm{f_S}^2.
\]
Take square roots and then use adjointness.
\end{proof}

\subsubsection{Wick input at a Hermite degree boundary}
For one sample define the row multiplier
\[
 Tf=y^{\mathsf T}f,\qquad T^*T=\mathcal M_Y
\]
as a quadratic-form identity on polynomials. All following estimates allow a
Hilbert-valued coefficient space for the other registers. H\"older is applied
only in the current sample variables, with that coefficient space kept intact.
Before the finite cutoff, identities involving multipliers and their adjoints
are understood as polynomial quadratic-form identities; after compression the
adjoints are ordinary finite-dimensional ones.

\begin{lemma}[Local degree-dependent row contraction]\label{c:lem:local}
Let $f_d$ have homogeneous local Hermite degree $d\ge1$. For every integer $p\ge2$,
\begin{equation}\label{c:eq:local}
 \boxed{\norm{Tf_d}_2^2\le
 M_p(r)\left(\frac{p+1}{p-1}\right)^d\norm{f_d}_2^2.}
\end{equation}
For a finite sum $f=\sum_{d\ge1}f_d$ with no local vacuum component, any
specified choices of integers $p_d\ge2$ give
\begin{equation}\label{c:eq:local-sum}
 \norm{Tf}_2^2\le3\sum_{d\ge1}
 M_{p_d}(r)\left(\frac{p_d+1}{p_d-1}\right)^d\norm{f_d}_2^2.
\end{equation}
\end{lemma}
\begin{proof}
\step{1. Apply H\"older on the original sample.}
The exponents are $2p$ and $2p/(p-1)$. Gaussian hypercontractivity
\eqref{eq:HC} therefore gives
\begin{align*}
 \norm{y^{\mathsf T}f_d}_2
 &\le\bigl\|\norm y\bigr\|_{L_{2p}}\norm{f_d}_{L_{2p/(p-1)}}\\
 &\le\bigl\|\norm y\bigr\|_{L_{2p}}
 \left(\frac{p+1}{p-1}\right)^{d/2}\norm{f_d}_2.
\end{align*}
Square and use \eqref{c:eq:wick-amgm} without relaxing $M_p(r)$. This proves
\eqref{c:eq:local}. The multiplier and the polynomial use the same Gaussian
sample; no independence between $y$ and $f_d$ is assumed.

\step{2. Sum homogeneous inputs using their output-degree overlap.}
Every coordinate of $y$ is pure second chaos: it is a linear combination of
$g_a h_b$ with coordinates in distinct Gaussian factors. Write $\mathcal F_e$
for exact local Hermite degree $e$, with the Hilbert coefficient space for the
other registers left intact; set $\mathcal F_e=\{0\}$ when $e<0$.
Applying \eqref{eq:Hermite} twice shows
\[
 Tf_d\in\mathcal F_{d-2}\oplus\mathcal F_d\oplus\mathcal F_{d+2}.
\]
At any output degree, at most three input degrees contribute.
Cauchy--Schwarz there, followed by orthogonality across output degrees, yields
$\norm{Tf}_2^2\le3\sum_d\norm{Tf_d}_2^2$. Apply \eqref{c:eq:local}
with the selected $p_d$ separately. This proves \eqref{c:eq:local-sum}.
\end{proof}

\paragraph{A simple explicit exponent policy.}
Choosing $p_d=2d+1$ in \eqref{c:eq:local} and using
$M_{2d+1}(r)\le(2d+1)(r+2d)$ gives
\begin{equation}\label{c:eq:local-simple}
 \norm{Tf_d}_2^2\le
 e\bigl[(2d+1)r+4d^2+2d\bigr]\norm{f_d}_2^2,
\end{equation}
since $(1+1/d)^d\le e$. The factor three in \eqref{c:eq:local-sum}
remains; the Wick improvement has not removed overlap or hypercontractivity.
Define the simple repeated-visit allowance
\begin{equation}\label{c:eq:Bq}
 \boxed{\widetilde B_q(r)=27e\,rq+192e\,q^2+24e\,q.}
\end{equation}
No free $\lambda$ is needed for this policy. Chain B retains its distinct $B_q(r;\lambda)$ in \eqref{eq:Bq};
it is not overwritten by $\widetilde B_q$.

\begin{lemma}[Sum of already-active sample interactions]\label{c:lem:repeated}
On the entire cutoff of \eqref{eq:cutoff},
\begin{equation}\label{c:eq:repeated}
 \norm{\mathcal R}\le\widetilde B_q(r).
\end{equation}
Its restriction to the vacuum sector is zero.
\end{lemma}
\begin{proof}
\step{1. Fix the exact active set and expose the centered form.}
On an active-set sector $S$, with $k=|S|\le q$,
\begin{equation}\label{c:eq:R-form}
 \ip f{\mathcal Rf}=\sum_{i\in S}\norm{T_if}_2^2-k\norm f_2^2.
\end{equation}
Indeed $Q_if=f$ for $i\in S$ and $Q_if=0$ otherwise, and
$\ip f{\mathcal X_if}=\norm{T_if}_2^2-\norm f_2^2$.
In particular $\mathcal R\succeq-kI$ on this sector.

\step{2. Sum each label's actual degree before relaxing any budget.}
Let $f=\sum_{\mathbf d}f_{\mathbf d}$ be the joint local-degree decomposition.
The number operators commute and their joint eigenspaces are orthogonal.
Using \eqref{c:eq:local-sum} with the policy \eqref{c:eq:local-simple},
\begin{align*}
 \sum_{i\in S}\norm{T_if}_2^2
 &\le3e\sum_{\mathbf d}
 \left[r\left(2\sum_{i\in S}d_i+k\right)
       +4\sum_{i\in S}d_i^2+2\sum_{i\in S}d_i\right]
       \norm{f_{\mathbf d}}_2^2\\
 &\le3e\bigl[9rq+64q^2+8q\bigr]\norm f_2^2
 =\widetilde B_q(r)\norm f_2^2.
\end{align*}
The only budget relaxations are
\[
 \sum_i d_i\le4q,\qquad k\le q,\qquad
 \sum_i d_i^2\le\left(\sum_i d_i\right)^2\le16q^2.
\]
Charging every active sample the total degree $4q$ separately would introduce
an unnecessary extra factor of $q$; summing the actual degrees avoids it.

\step{3. Control both sides and preserve the cutoff.}
Equation~\eqref{c:eq:R-form} gives
$-kI\preceq\mathcal R\preceq\widetilde B_q I$.
Since $\widetilde B_q\ge q\ge k$, its norm is at most $\widetilde B_q$.
These form inequalities survive the finite cutoff. Sum over orthogonal
exact-support sectors; on $S=\varnothing$ the operator vanishes.
\end{proof}

\subsection{C4. Propagate the hybrid budgets and recover the endpoint}
Put $K=\min(q,m)$ and retain the Wick activation coefficients
\begin{equation}\label{c:eq:off}
 a_k^{(C)}=\frac{\sqrt{(k+1)(m-k)v_r^{(C)}}}{m},
 \qquad v_r^{(C)}=4r+4,\quad 0\le k<K.
\end{equation}
Let $G_{C,q}$ be the nonnegative symmetric tridiagonal matrix with
$(G_{C,q})_{00}=0$, nonvacuum diagonal $\widetilde B_q(r)/m$, and off-diagonal
$a_k^{(C)}$. Activation, return, and repeated-visit estimates give
$n(M_qf)\le G_{C,q}n(f)$, so the vacuum identity yields
\begin{equation}\label{c:eq:stationary}
 \mathbb E\tr Z^{2q}\le r\|G_{C,q}^{q}e_0\|_2^2.
\end{equation}
Here $n(f)$ is the vector of active-count sector norms from
\eqref{eq:sector-step}; it does not encode an independent replacement model.

\begin{theorem}[Chain C closed-form moment and tail]\label{c:thm:closed}
For every integer $q\ge1$, with $c_q=((2q-1)!!)^{1/(2q)}$ and
$\eta_q=(r/\delta)^{1/(2q)}$,
\begin{align}
 (\mathbb E\tr Z^{2q})^{1/(2q)}
 &\le r^{1/(2q)}\left[c_q\sqrt{\frac{4r+4}{m}}+
                 \frac{\widetilde B_q(r)}{m}\right],\label{c:eq:closed}\\
 \mathbb P\!\left\{\|Z\|>
 \eta_q\left[c_q\sqrt{\frac{4r+4}{m}}+
                 \frac{\widetilde B_q(r)}{m}\right]\right\}&\le\delta.
 \label{c:eq:tail}
\end{align}
Consequently Chain C also proves the common endpoint \eqref{eq:main-size}.
\end{theorem}
\begin{proof}
\step{1. Preserve the comparison vacuum.}
As in the fully proved comparison in Theorem~\ref{thm:B-closed},
$G_{C,q}$ is entrywise dominated by
$bI+aJ_q$, with $b=\widetilde B_q/m$ and $a=\sqrt{(4r+4)/m}$.
The positive iteration preserves this domination. The scalar Hermite recurrence
identifies $\|(bI+aJ_q)^qe_0\|^2$ with $\mathbb E(b+a\xi)^{2q}$,
where $\xi$ is standard Gaussian. This comparison random variable changes no
original Gaussian sample.

\step{2. Take the moment root and apply Markov.}
Scalar Minkowski gives the first inequality. The spectral inequality
$\|Z\|^{2q}\le\tr Z^{2q}$ and Markov give the second. With
$q=\lceil\log(r/\delta)\rceil$, $\eta_q\le\sqrt e$ and
$c_q\le\sqrt{2q}$, the allowance has order
$\sqrt{rq/m}+rq/m+q^2/m+q/m$. Since $r,q\ge1$ and $\varepsilon<1$,
requiring $m\ge C(\varepsilon^{-2}rt+\varepsilon^{-1}t^2)$, with
$t=\log(4r/\delta)$, controls all terms. This is the same rate as A and B.
\end{proof}

\paragraph{Stage-dependent hybrid propagation.}
For $1\le j\le q$, let $G_C^{[j]}$ have the same off-diagonal
$a_k^{(C)}$, zero vacuum diagonal, and nonvacuum diagonal
$\widetilde B_j(r)/m$. The exact prefix after $j-1$ steps belongs to the cutoff
at total degree $4(j-1)$; its next multiplication needs only degree $4j$.
Apply the C-contractions at that next cutoff and iterate sector norms in order:
\begin{equation}\label{c:eq:simple-stage}
 \mathbb E\tr Z^{2q}
 \le r\|G_C^{[q]}\cdots G_C^{[1]}e_0\|_2^2
 \le r\|G_{C,q}^{q}e_0\|_2^2.
\end{equation}
The proof is the same positive iteration as Theorem~\ref{thm:stagewise}, with
\emph{the C-budgets} supplied by Lemmas~\ref{c:lem:fresh} and
\ref{c:lem:repeated}; the operator identity is unchanged. The last inequality
uses $\widetilde B_j\le\widetilde B_q$. No stage matrices are commuted.

\subsection{Refinement of C4: exact radial products and degree allocations}
\label{c:sec:refined}
The simple C-stagewise bound \eqref{c:eq:simple-stage} still uses AM--GM for $M_p(r)$, replaces each
$(1+1/d)^d$ by $e$, and uses one diagonal allowance for every active count.
We now retain those quantities. This is a finite deterministic refinement of
the same operator proof, not a new random sketch or a simulation.

Fix the terminal $q$ and the finite exponent menu
\[
 \mathcal P_q=\{2,3,\ldots,8q+1\}.
\]
For every even $d\in\{2,4,\ldots,4q\}$ choose a fixed $p_d\in\mathcal P_q$ and put
\begin{equation}\label{c:eq:wick-degree-cost}
 h_d(r)=M_{p_d}(r)\left(\frac{p_d+1}{p_d-1}\right)^d,
 \qquad b_d(r)=3h_d(r)-1.
\end{equation}
Any fixed choices give valid local upper bounds. For the additional comparison
with the simple certificate, we require
\begin{equation}\label{c:eq:choice-policy}
 h_d(r)\le M_{2d+1}(r)\left(\frac{d+1}{d}\right)^d.
\end{equation}
For instance one may minimize over $\mathcal P_q$, since the menu contains
$2d+1$. Alternatively a predetermined list satisfying \eqref{c:eq:choice-policy}
is sufficient; no proof that the list minimizes the menu is needed. The lists
used in the numerical table are specified and checked in Appendix~\ref{app:checks}.

For $1\le k\le j\le q$, define
\begin{equation}\label{c:eq:wick-budget}
 \mathfrak D_{j,k}(r)=
 \max_{\substack{d_1,\ldots,d_k\in2\mathbb N\\d_1+\cdots+d_k\le4j}}
       \sum_{i=1}^k b_{d_i}(r),\qquad \mathfrak D_{j,0}=0.
\end{equation}
The symbol $\mathfrak D$ denotes this degree maximum, not the radial product
$D_{r,p}$ in \eqref{c:eq:wick-parameters}. All the costs use the same terminal
choice of $(p_d)$ throughout the $q$ stages.

\begin{lemma}[Exact-degree repeated-visit budget]\label{c:lem:wick-budget}
On the stage-$j$ cutoff restricted to $k$ active samples,
\[
 \norm{\mathcal R|_{\mathcal K_k}}\le\mathfrak D_{j,k}(r).
\]
For every admissible exponent choice $\mathfrak D_{j,k}\ge k$.
If \eqref{c:eq:choice-policy} also holds, then
$\mathfrak D_{j,k}\le\widetilde B_j(r)$.
\end{lemma}
\begin{proof}
\step{1. Keep the centering term in each local cost.}
On a fixed active set, apply \eqref{c:eq:local-sum} in \eqref{c:eq:R-form}
and then the joint degree decomposition. For every component the upper cost is
$\sum_i(3h_{d_i}-1)=\sum_i b_{d_i}$; the $-1$ terms sum to the exact $-k$
centering cost. The original symmetries require even $d_i\ge2$ and the next
cutoff requires total degree at most $4j$. This proves the stated upper form bound.

\step{2. Bound the negative side and compare policies.}
The negative form bound is $-kI$. Since $M_p(r)\ge r$ and the other factor
in $h_d$ is at least one, $h_d\ge r\ge1$ and hence $b_d\ge2$.
Thus $\mathfrak D_{j,k}\ge k$, controlling both signs. Under
\eqref{c:eq:choice-policy}, each $h_d$ is at most the simple local allowance
in \eqref{c:eq:local-simple}. Sum as in Lemma~\ref{c:lem:repeated} to get
$\mathfrak D_{j,k}\le\widetilde B_j$.
\end{proof}

Let $K=\min(q,m)$ and $a_k=\sqrt{(k+1)(m-k)v_r^{(C)}}/m$. Define
$\widehat G_C^{[j]}$ on $0,\ldots,K$ by
\begin{align}
 (\widehat G_C^{[j]})_{00}&=0,\nonumber\\
 (\widehat G_C^{[j]})_{kk}&=
 \begin{cases}\mathfrak D_{j,k}(r)/m,&1\le k\le\min(j,K),\\
 0,&j<k\le K,\end{cases}\nonumber\\
 (\widehat G_C^{[j]})_{k+1,k}&=(\widehat G_C^{[j]})_{k,k+1}=a_k.
 \label{c:eq:wick-stage-matrix}
\end{align}
The zero entries beyond the currently reachable active counts are padding.

\begin{theorem}[Wick radial/degree vacuum certificate]\label{c:thm:wick-refined}
For the original independent Gaussian law and every fixed admissible exponent list,
\begin{equation}\label{c:eq:wick-stage-moment}
 \boxed{\E\tr Z^{2q}\le
 r\norm{\widehat G_C^{[q]}\widehat G_C^{[q-1]}\cdots\widehat G_C^{[1]}e_0}_2^2.}
\end{equation}
In particular,
\begin{equation}\label{c:eq:wick-stage-tail}
 r\norm{\widehat G_C^{[q]}\cdots\widehat G_C^{[1]}e_0}_2^2
 \le\delta\varepsilon^{2q}
\end{equation}
is a sufficient original-law embedding certificate. Under \eqref{c:eq:choice-policy}
it is no worse than the simple stage-dependent certificate at the same $q,r,m$.
\end{theorem}
\begin{proof}
For the exact prefixes $f_j=\mathcal M_Z^j(e_a\otimes\Omega)$, the next cutoff
identity is still $f_j=M_jf_{j-1}$ as in Theorem~\ref{thm:stagewise}.
Use Lemma~\ref{c:lem:fresh} for activation and return, and
Lemma~\ref{c:lem:wick-budget} for the diagonal interaction. Hence
$n(f_j)\le\widehat G_C^{[j]}n(f_{j-1})$ componentwise. Induct in stage order,
square the terminal sector norm, and sum the $r$ external vacuum seeds.
This proves \eqref{c:eq:wick-stage-moment}; Markov gives
\eqref{c:eq:wick-stage-tail}. Under \eqref{c:eq:choice-policy},
$\widehat G_C^{[j]}\le G_C^{[j]}$ entrywise, so positivity gives the claimed
comparison. No commutation of the stage matrices is used.
\end{proof}

\subsection{Completion, scope, and finite-moment transfer}
Either the simple or the refined C-certificate, followed by Markov, supplies
the same arbitrary-fixed-subspace embedding guarantee. The common lower bound
in Section~\ref{sec:lower} therefore matches C as well as A and B at fixed
confidence and distortion. The constant $4r+4$ is a uniform covariance envelope,
not a sharp-covariance claim; the radial envelope is attained, but subsequent
H\"older, degree-overlap, and scalar-majorant estimates can lose information.

The $2q$-wise sample-block transfer of Corollary~\ref{cor:limited} applies to C:
it transfers the \emph{untruncated} polynomial moment
$\mathbb E\tr Z^{2q}$. It neither constructs a finite-seed Gaussian generator nor
transfers an inverse bound or the common lower-tail obstruction. As in B,
no positive-coefficient cone, general descent theorem, or covariance-spectrum
renewal theorem is needed here. The hybrid uses exact contractions at the
local coefficient boundary, then the full-space Hermite framework.
\section{A common matching lower bound}\label{sec:lower}
All three upper proofs concern the same distribution, so they share one lower-bound
argument. A single large sample suffices; no spectral independence of matrix
entries is required.

\begin{theorem}[Common-factor obstruction]\label{thm:lower}
Take $n_1=r$, $n_2=2$, and $Ue_a=e_a\otimes e_1$ for $1\le a\le r$.
For all sufficiently large $r$, uniformly over integers
$1\le m\le\tfrac14r\log r$,
\[
 \Prob\{\norm Z>1/2\}\longrightarrow1\qquad(r\longrightarrow\infty).
\]
\end{theorem}
\begin{proof}
For this subspace,
\[
 y_i=\eta_i\xi_i,\qquad \eta_i\sim N(0,1),\quad\xi_i\sim N(0,I_r),
\]
with all scalar and vector variables independent.

\step{1. Below $r$ samples, rank already forces failure.}
If $m<r$, the sample covariance $m^{-1}\sum_i y_iy_i^{\mathsf T}$ has rank at
most $m$, so it has a zero eigenvalue. Therefore $\norm Z\ge1$ deterministically.

\step{2. Select the large scalar without biasing its vector factor.}
Suppose $m\ge r$ and let $J$ be the smallest index maximizing $\eta_i^2$.
The index depends only on the scalars. Conditional on all of them, $\xi_J$ still
has law $N(0,I_r)$; thus it is independent of those scalars.
Positivity of the remaining covariance summands gives
\begin{equation}\label{eq:lower-spike}
 \lambda_{\max}\left(\frac1m\sum_i y_iy_i^{\mathsf T}\right)
 \ge\frac{\eta_J^2\norm{\xi_J}^2}{m}.
\end{equation}

\step{3. Lower-bound the selected scalar and the vector length.}
For $m\ge3$, integrate the Gaussian density on
$[a,a+1/a]$ and its reflection, where $a=\sqrt{\log m}\ge1$.
Since $(a+1/a)^2\le a^2+3$, this gives an absolute $c_1>0$ such that
\[
 p_m:=\Prob\{\eta^2\ge\log m\}
 \ge c_1\frac{m^{-1/2}}{\sqrt{\log m}}.
\]
Independence implies
\[
 \Prob\{\eta_J^2<\log m\}=(1-p_m)^m
 \le\exp\left(-c_1\frac{\sqrt m}{\sqrt{\log m}}\right).
\]
Also $\E e^{-\norm\xi^2/2}=2^{-r/2}$, so exponential Markov gives
\[
 \Prob\{\norm{\xi_J}^2<r/2\}
 \le e^{r/4}2^{-r/2}\le e^{-r/16}.
\]
Consequently, with probability at least
\begin{equation}\label{eq:lower-prob}
 1-\exp\left(-c_1\frac{\sqrt m}{\sqrt{\log m}}\right)-e^{-r/16},
\end{equation}
the right side of \eqref{eq:lower-spike} is at least $r\log m/(2m)$.

\step{4. Read off the necessary sample order.}
If $r\le m\le\tfrac14r\log r$, then $\log m\ge\log r$, and hence
$r\log m/(2m)\ge2$. On the event above, $\norm Z\ge1>1/2$.
The probability in \eqref{eq:lower-prob} tends to one uniformly throughout this
range, since $m\ge r\to\infty$. Together with the rank argument this proves the
claim.
\end{proof}
At $\varepsilon=1/2$, $\delta=0.01$, one has
$t=\log(400r)=O(\log(er))$ and $t^2=O(r\log(er))$.
Any of the three upper proofs and Theorem~\ref{thm:lower} therefore establish the
$\Theta(r\log(er))$ assertion of Theorem~\ref{thm:common-endpoint}.
The obstruction is worst-case: a particular subspace may need fewer rows.

\Needspace{10\baselineskip}
\section{Comparing A, B, and C: endpoints and constants}\label{sec:comparison}
\subsection{Fix the labels before comparing a number}
Chain A means the concentration proof of Section~\ref{sec:A}, including its
original raw proxy $9r$ and moment scale $72(r+s)$. Chain B means the framework
proof of Section~\ref{sec:B}, including $v_r=9r-1$ and $B_j(r;\lambda)$.
Chain C means the Wick--Hermite proof of Section~\ref{sec:C}, with
$v_r^{(C)}=4r+4$, $\widetilde B_j$, and optionally $\mathfrak D_{j,k}$.
\emph{None of these labels changes meaning when a sharper estimate is discussed.}

For the principal comparison, define
\begin{align}
 U_A&=\sqrt{\frac{18rt}{m}}+\frac{144(r+s)t}{m}+3e^{-s},
 \quad t=\log(4r/\delta),\quad s=\lceil\log(4m/\delta)\rceil;
 \label{eq:UA}\\
 U_B&=\left[\frac r\delta\|G^{[q]}\cdots G^{[1]}e_0\|_2^2\right]^{1/(2q)};
 \label{eq:UBdynamic}\\
 U_C&=\left[\frac r\delta\|\widehat G_C^{[q]}\cdots\widehat G_C^{[1]}e_0\|_2^2\right]^{1/(2q)}.
 \label{eq:UCdynamic}
\end{align}
In the cover diagram $G_B^{[j]}$ is simply an explicit alias for B's
$G^{[j]}$ in \eqref{eq:stage-matrix}. The C matrices in \eqref{c:eq:wick-stage-matrix}
are different majorants of the same underlying multiplication. Each displayed
allowance being at most $\varepsilon$ certifies $\mathbb P\{\|Z\|>\varepsilon\}\le\delta$.

\subsection{The local constants, side by side}
\begin{table}[H]\centering\small
\setlength{\tabcolsep}{4pt}\renewcommand{\arraystretch}{1.18}
\begin{tabularx}{\textwidth}{@{}P{.22\textwidth}YYY@{}}
\toprule
\textbf{Quantity}&\textbf{A: concentration}&\textbf{B: framework}&\textbf{C: Wick--Hermite}\\
\midrule
Row input at $L_{2p}$, squared&
$r(4p-3)+(2p-2)^2$&Same baseline&
$M_p(r)=(\kappa_pD_{r,p})^{1/p}\le p(r+p-1)$\\[3pt]
Raw second moment&$9r$&$9r$&$4r+5$\\[3pt]
Centered activation variance&Not the assembly parameter&$9r-1$&$4r+4$\\[3pt]
Exponential-moment scale&$72(r+s)$&Not used&Not used\\[3pt]
Simple repeat allowance&Not used&
$99e\,rq+768e\,q^2$ at $\lambda=2$&
$27e\,rq+192e\,q^2+24e\,q$\\[3pt]
Stronger propagation&Matrix Laplace&
Ordered $G^{[q]}\cdots G^{[1]}$&
Ordered $\widehat G_C^{[q]}\cdots\widehat G_C^{[1]}$; active-degree maximum\\[3pt]
Untruncated trace moments&Not supplied by its displayed tail certificate&Yes&Yes\\
\bottomrule
\end{tabularx}
\caption{The fixed A/B baselines and the separate C refinement. The numerical
inequalities are proved in their own sections; shared symbols do not overwrite
old constants.}\label{tab:three-constants}
\end{table}

For B also retain the closed and stationary allowances
\begin{align}
 U_{B,\mathrm{cl}}&=\eta_q\left[c_q\sqrt{\frac{9r-1}{m}}+
 \frac{B_q(r;\lambda)}m\right],\label{eq:UBclosed}\\
 U_{B,\mathrm{st}}&=\left[\frac r\delta\|G_q^qe_0\|_2^2\right]^{1/(2q)}.
 \label{eq:UBstatic}
\end{align}
At identical $m,r,q,\lambda,\delta$,
$U_B\le U_{B,\mathrm{st}}\le U_{B,\mathrm{cl}}$.
For C denote the allowance from \eqref{c:eq:simple-stage} by
$U_{C,\mathrm{simple}}$. Under \eqref{c:eq:choice-policy},
$U_C\le U_{C,\mathrm{simple}}$ at identical parameters. For the comparison
$\lambda=2$, $\widetilde B_j\le B_j(r;2)$ because their difference is
$72e\,rj+576e\,j^2-24e\,j>0$. Also $4r+4\le9r-1$ for $r\ge1$.
Thus $U_C\le U_{C,\mathrm{simple}}\le U_B$ at the same parameters and that
baseline policy. The $r=10$ table uses $\lambda=1$; there the analogous
budget inequality is checked for the displayed stages. No universal comparison
with every optimized concentration proof is inferred.

\subsection{The main numerical comparison, with proof parameters}
Table~\ref{tab:three-main-counts} is the principal A/B/C comparison.
The model and $\varepsilon=1/2,\delta=0.01$ are identical. B and C use the same
terminal $q$, so the reduction from B to C is not obtained by changing the
moment order. B uses $\lambda=1$ at $r=10$ and $\lambda=2$ otherwise.
C's exact exponent lists appear in Table~\ref{tab:wick-exponents}.

\begin{table}[H]\centering\small
\setlength{\tabcolsep}{7pt}\renewcommand{\arraystretch}{1.2}
\begin{tabular}{@{}rrrrrr@{}}
\toprule
$r$&$q$&$\lambda_B$&\textbf{A}&\textbf{B}&\textbf{C}\\
\midrule
10&3&1&91\,000&60\,000&12\,800\\
100&5&2&580\,000&472\,000&111\,000\\
1\,000&7&2&6\,200\,000&5\,150\,000&1\,250\,000\\
10\,000&8&2&72\,000\,000&61\,000\,000&15\,000\,000\\
\bottomrule
\end{tabular}
\caption{All parameters of the primary A/B/C certificates; C uses the
specified radial/degree policy, not an unspecified optimization.}
\label{tab:main-parameters}
\end{table}

The comparison reduces the displayed B row counts by factors approximately
$4.69$, $4.25$, $4.12$, and $4.07$, respectively. These ratios concern the
selected certificates, not the true minimum sample complexity. In particular,
A and B are intentionally the original baselines: their separate identity is
what makes the contribution of C visible.

\subsection{The gains and losses behind the constants}
Let $h=\log(r/\delta)$ and $q=\lceil h\rceil$, with $h\to\infty$.
Stirling gives $c_q\sim\sqrt{2q/e}$ and $\eta_q\sim\sqrt e$.
Consequently the fresh-sample term has asymptotic coefficient
\[
 \text{A and B:}\quad\sqrt{18rh/m},\qquad
 \text{C:}\quad\sqrt{8rh/m},
\]
when also $r\to\infty$. The improved covariance envelope alone reduces
that coefficient by a factor $3/2$.
At $\lambda=2$ and $r/h\to\infty$, the simple repeated-visit coefficient
in the closed form changes from $99e^{3/2}\approx443.69$ for B to
$27e^{3/2}\approx121.01$ for C; the exact-degree certificate can improve
further. Neither coefficient is a globally optimal constant.

Retaining the vacuum rather than replacing it by Schur's full-norm bound
still avoids the asymptotic fresh-term factor $\sqrt{2e}$ in either framework
route. C has \emph{not} removed the three-degree-overlap factor, hypercontractivity
on arbitrary polynomial coefficients, or the block-norm majorization itself.
Its sharp radial bound does not imply an optimal final embedding constant.

\begin{takeaway}[title={Same endpoint, a distinguishable hybrid}]
A, B, and C all yield
$O(\varepsilon^{-2}r\log(4r/\delta)+\varepsilon^{-1}\log^2(4r/\delta))$,
and the common fixed-confidence lower bound gives $\Theta(r\log(er))$.
The difference is the proved coefficient budget and the resulting explicit
certificate. C has a separate derivation and improves the stated A/B baselines.
Appendix~\ref{app:A-wick} records a fair additional comparison after applying
Wick inputs to A as well; it prevents interpreting this table as intrinsic
superiority of one general method.
\end{takeaway}
\section{Scope, provenance, and verification status}\label{sec:scope}
The factors throughout are independent real standard Gaussian vectors, and the
subspace is deterministic. No step establishes the same improvement for arbitrary
isotropic sub-Gaussian factor vectors, higher tensor orders, or subspaces chosen
after observing the sketch. The lower bound is worst-case, not a classification
of every subspace by Schmidt geometry. Optimal joint $\varepsilon,\delta$
dependence is not claimed.

Chain A reflects the coefficient-retention viewpoint but is a standalone
concentration argument. Chains B and C use the framework's exact multiplication,
graded-domain reduction, coefficient-sensitive block estimates, and vacuum/energy
propagation. Neither needs the positive-coefficient cone, the graph
connected-creation theorem, general encoding/descent tests, signed conditional
cores, or covariance-spectrum renewal. In particular, their Gaussian
hypercontractivity bounds apply to arbitrary polynomial coefficients.

The present edition corrects the preceding integration's organization: it
restores A and B as the original baselines and gives the Wick--Hermite
refinement its own Chain C. The supplied three-chain rMPS manuscript is a
layout reference only. All theorem statements and numerical quantities here
concern the same Khatri--Rao embedding model. The prior operator-domain
correction and expanded proof exposition are retained.

The starting material was the two proof notes developed in this conversation:
\emph{Sharp fixed-confidence order-two Gaussian Khatri--Rao embeddings} and
\emph{Gaussian Khatri--Rao embeddings through the graded multiplication pipeline},
both dated 16 September 2026. This document preserves their baseline application arguments and integrates
the Wick supplement as the separate Chain C~\cite{WickNote,RMPS,Janson}.
The organizing framework is the supplied v11 manuscript~\cite{Framework}.
The Gaussian inequalities, moment-estimation methods, and matrix-analysis inputs
are classical and are credited at their points of use. Local proof provenance
in a conversation is not evidence of publication priority. In particular, the
row-moment order has the prior implication explained after Lemma~\ref{lem:row},
and Chain A's general proof strategy has the predecessors recorded in
Section~\ref{sec:attribution-intro}.

\clearpage
\subsection{Attribution by ingredient}\label{sec:provenance-table}
\begingroup\small\setlength{\tabcolsep}{5pt}\renewcommand{\arraystretch}{1.2}
\begin{longtable}{@{}P{.23\textwidth}P{.36\textwidth}P{.34\textwidth}@{}}
\toprule
\textbf{Ingredient}&\textbf{Prior source or standard step}&\textbf{Role of the present note}\\
\midrule
Sketch construction and earlier guarantees &
Tensor random projection~\cite{Sun}; fixed-vector sketching~\cite{Ahle};
rank-one probing~\cite{BK}; earlier OSE analysis~\cite{BGKL,BM}. &
Same sketch, with the precise real-Gaussian, order-two theorem proved here.
No claim to introduce the sketch.\\
Gaussian functional inequalities &
Nelson hypercontractivity; Gross logarithmic Sobolev theory;
Aida--Stroock moment methods~\cite{Nelson,Gross,GrossSurvey,AidaStroock}. &
Re-derivations fix constants and the Hilbert-valued normalization.\\
Projected-row moments &
The order follows from the Hilbert case of~\cite[Corollary~12]{ALM}.
Fourth moments are Isserlis calculations~\cite{Isserlis}. &
Explicit baseline row bound and its use in A/B; not a new general chaos-moment
order.\\
Concentration chain &
Khatri--Rao truncation strategy~\cite{BM}, with the version-sensitive
predecessor~\cite{SVBv1,SVBv2}; Lieb--Tropp matrix Laplace/Bernstein
\cite{Lieb,Tropp}. &
Rank-one matrix-moment bookkeeping, all-order cutoff calculation, bias control,
and one explicit certificate.\\
Graded representation &
Classical chaos, probability-law representations, and Jacobi moments
\cite{Wiener,AccardiBozejko,GolubWelsch}; organization in~\cite{Framework}. &
Exact support/degree restriction for this model. Representation alone is not
claimed as novel.\\
Coefficient and energy estimates &
Coefficient-sensitive chaos methods~\cite{ALM,BLNvH}; block norms and positive
majorization; energy interface in~\cite{Framework}. &
Specific activation, local-degree repeated-visit, and ordered vacuum
calculations with displayed constants.\\
Wick--Hermite hybrid &
Classical Gaussian pairings and chaos~\cite{Isserlis,Janson}; radialization
recorded in~\cite{RMPS}; supplied calculations in~\cite{WickNote}. &
Sharp radial product, orthogonal mixed-fourth estimate, and refined C-stage
budgets proved here. Internal provenance is not external priority.\\
Lower bound and moment transfer &
Elementary positivity, Gaussian tails, and low-degree moment matching. &
Explicit common-factor witness and sample-block transfer proof; no claim to
invent these general proof techniques or a finite-seed construction.\\
\bottomrule
\end{longtable}
\endgroup

The numerical comparison is internal: all constants come from the implementations
proved here, not from an optimization over the cited authors' theorems. Showing
that one displayed certificate is smaller than another establishes no universal
ranking of the methods. The source checks identify relevant antecedents but are
not an exhaustive literature search establishing originality of the final
embedding theorem.

\subsection{Verification status}

The accompanying checks verify exact finite scalar inequalities, selected
high-precision numerical certificates, and normalization benchmarks. They do not
constitute an independent proof audit, a Lean certificate, a test of every
possible parameter value, or a claim of publication priority. The general
statements are justified by the analytic proofs, not by the numerical tests.


\clearpage
\appendix
\section{Gaussian inputs and Hermite normalization}\label{app:gaussian}
This appendix records short standard semigroup arguments for the Gaussian
inequalities used in Section~\ref{sec:common}. They are included to fix constants,
not asserted as new results; see Nelson~\cite{Nelson} and
Gross~\cite{Gross,GrossSurvey}.

\subsection{Ornstein--Uhlenbeck semigroup and logarithmic Sobolev inequality}
For a smooth function $F$ on $\R^N$, let
\[
 (\mathsf P_tF)(x)=\E F\bigl(e^{-t}x+\sqrt{1-e^{-2t}}\,G\bigr),
 \qquad G\sim N(0,I_N).
\]
Standard Gaussian measure is invariant under $\mathsf P_t$; this follows directly
because a sum of independent centered Gaussian vectors with squared coefficients
summing to one is standard Gaussian. Differentiation of the kernel gives the
generator $\mathscr L=\Delta-x\cdot\nabla$, and Gaussian integration by parts gives
\[
 \E f\mathscr Lg=-\E\ip{\nabla f}{\nabla g},\qquad
 \nabla\mathsf P_tF=e^{-t}\mathsf P_t\nabla F.
\]

\step{1. Differentiate the entropy along the semigroup.}
Assume first that $F$ is smooth, bounded, and bounded below by a positive constant.
Set $u_t=\mathsf P_tF$. Since $\E u_t=\E F$,
\[
 \frac{d}{dt}\Ent(u_t)
 =\E(\mathscr Lu_t)\log u_t
 =-\E\frac{\norm{\nabla u_t}^2}{u_t}.
\]
As $t\to\infty$, $u_t\to\E F$ and its entropy tends to zero, so integrating this
derivative recovers $\Ent(F)$.

\step{2. Bound the dissipation pointwise.}
The gradient commutation identity and weighted Cauchy--Schwarz give pointwise
\[
 \frac{\norm{\nabla\mathsf P_tF}^2}{\mathsf P_tF}
 \le e^{-2t}\mathsf P_t\left(\frac{\norm{\nabla F}^2}{F}\right).
\]

\step{3. Integrate in space, then in time.}
Integrate in Gaussian measure, use invariance, and then integrate over $t$:
\[
 \Ent(F)
 \le\int_0^\infty e^{-2t}\,dt\;
 \E\frac{\norm{\nabla F}^2}{F}
 =\frac12\E\frac{\norm{\nabla F}^2}{F}.
\]
Taking $F=f^2$ gives
\[
 \Ent(f^2)\le2\E\norm{\nabla f}^2,
\]
which is \eqref{eq:LSI}. Positive regularization, truncation, and smooth
approximation extend the inequality to the functions with finite energy used
in the main proof.

\subsection{Hypercontractivity and homogeneous chaos}
\step{1. A monotone norm along the semigroup.}
Let $f>0$ be smooth and bounded, $u_t=\mathsf P_tf$, and
$p(t)=1+(p_0-1)e^{2t}$ with $p_0>1$. Differentiation yields
\begin{align*}
 \frac{d}{dt}\log\norm{u_t}_{L_{p(t)}}
 &=\frac{p'(t)}{p(t)^2}
   \frac{\Ent(u_t^{p(t)})}{\E u_t^{p(t)}}
 -(p(t)-1)\frac{\E u_t^{p(t)-2}\norm{\nabla u_t}^2}{\E u_t^{p(t)}}.
\end{align*}
Apply \eqref{eq:LSI} to $u_t^{p(t)/2}$:
\[
 \Ent(u_t^{p(t)})
 \le\frac{p(t)^2}{2}\E u_t^{p(t)-2}\norm{\nabla u_t}^2.
\]
Since $p'(t)=2(p(t)-1)$, the derivative of the logarithmic norm is nonpositive.
Thus $\norm{\mathsf P_tf}_{L_{p(t)}}\le\norm f_{L_{p_0}}$.
Positivity of the semigroup and $|\mathsf P_tf|\le\mathsf P_t|f|$ extend this to
signed functions, then approximation extends it to the relevant $L_p$ spaces.

\step{2. Specialize to homogeneous chaos.}
The Hermite generating function
\[
 e^{zx-z^2/2}=\sum_{\ell=0}^\infty\operatorname{He}_\ell(x)\frac{z^\ell}{\ell!}
\]
shows that $\mathsf P_t$ acts on homogeneous total degree $d$ by $e^{-dt}$.
Take $p_0=2$ and $u=1+e^{2t}$. For a scalar homogeneous Hermite polynomial $F_d$,
\[
 e^{-dt}\norm{F_d}_{L_u}\le\norm{F_d}_{L_2},
 \quad\text{hence}\quad
 \norm{F_d}_{L_u}\le(u-1)^{d/2}\norm{F_d}_{L_2}.
\]

\step{3. Allow Hilbert-space coefficients.}
For a Hilbert-valued polynomial with orthonormal coordinates $F_{d,a}$,
Minkowski in $L_{u/2}$ gives
\[
 \norm{F_d}_{L_u}^2
 =\left\|\sum_a|F_{d,a}|^2\right\|_{L_{u/2}}
 \le\sum_a\norm{F_{d,a}}_{L_u}^2
 \le(u-1)^d\sum_a\norm{F_{d,a}}_{L_2}^2.
\]
Finite-dimensional projections and monotone convergence justify arbitrary
Hilbert coefficient spaces. This proves \eqref{eq:HC} with the stated constant.

\subsection{The scalar Jacobi comparison is exact through the needed moment}
For $h_\ell=\operatorname{He}_\ell/\sqrt{\ell!}$, the generating function also
yields \eqref{eq:Hermite}. The first $q+1$ Hermite functions therefore represent
compressed scalar multiplication by $\xi\sim N(0,1)$ with Jacobi matrix $J_q$.
An induction on the prefix length gives
\[
 (bI+aJ_q)^je_0\quad\longleftrightarrow\quad(b+a\xi)^j,
 \qquad 0\le j\le q.
\]
The Gaussian moment recurrence, obtained by integration by parts, is
$\E\xi^{2q}=(2q-1)\E\xi^{2q-2}$. It follows that
\[
 \norm{J_q^qe_0}^2=(2q-1)!!,\qquad
 \norm{(bI+aJ_q)^qe_0}^2=\E(b+a\xi)^{2q}.
\]
These are vacuum identities, not traces of the entire Jacobi matrix. They are
the Hermite case of the classical Jacobi/Gaussian-quadrature correspondence
\cite{GolubWelsch}; the Gaussian moment recurrence is also the one-variable case
of the product-moment formula~\cite{Isserlis}.

\section{An all-parameter verification of the constant 72}\label{app:72}
We prove \eqref{eq:72-scalar} without optimizing the constant. Set
\[
 a_p=8p-11,\qquad b_p=(4p-6)^2,\qquad
 F_p(r)=\frac{(a_pr+b_p)^{p-1}}{r(r+p)^{p-2}}.
\]
For integer $p\ge3$, the sign of the logarithmic derivative in $r$ is the sign of
\[
 \bigl[-8p^3+37p^2-62p+36\bigr]r-b_pp.
\]
Indeed, multiply $(p-1)a_p/(a_pr+b_p)-1/r-(p-2)/(r+p)$ by its positive
common denominator and collect terms. If $z=p-3\ge0$, the bracket becomes
\[
 -8z^3-35z^2-56z-33<0.
\]
Thus $F_p(r)$ decreases for $r\ge1$, so it is enough to prove the inequality at
$r=1$. Since $a_p+b_p=(4p-5)^2$, the ratio of its left to right sides is
\begin{equation}\label{eq:Rp}
 \rho_p=\frac{2(4p-5)^{2p-2}}
 {3p!\,72^{p-2}(p+1)^{p-2}}.
\end{equation}

\step{Finite orders $3\le p\le10$.}
The following rational upper bounds are verified by integer cross-multiplication
in \eqref{eq:Rp}; each is less than one.
\begin{center}
\footnotesize\renewcommand{\arraystretch}{1.25}\setlength{\tabcolsep}{5pt}
\begin{tabular}{@{}rrrrrrrrr@{}}
\toprule
$p$&3&4&5&6&7&8&9&10\\
\midrule
Upper bound on $\rho_p$&$927/1000$&$380/1000$&$177/1000$&$88/1000$&$46/1000$&$25/1000$&$14/1000$&$8/1000$\\
\bottomrule
\end{tabular}
\end{center}
For example $\rho_3=2401/2592<927/1000$.
All eight exact fractions are also recorded by the companion verification script.

\step{All remaining orders $p\ge11$.}
The elementary inequality $p!\ge(p/e)^p$ follows from
$\sum_{k=1}^p\log k\ge\int_1^p\log x\,dx$.
Using $4p-5\le4p$ and $p+1\ge p$ in \eqref{eq:Rp} gives
\[
 \rho_p\le\frac{32e^2}{3}\left(\frac{2e}{9}\right)^{p-2}.
\]
Since $e<11/4$ and $11/18<1$, for every $p\ge11$,
\[
 \rho_p<\frac{242}{3}\left(\frac{11}{18}\right)^9
 =\frac{285311670611}{297538935552}<1.
\]
This proves \eqref{eq:72-scalar} for every real $r\ge1$ and integer $p\ge3$.
It is an all-parameter analytic certificate with only eight finite integer checks,
not a numerical sweep over dimensions or high moment orders.

\section{Numerical certificates and reproducibility}\label{app:checks}
\subsection{One recurrence, with explicit chain-dependent coefficients}
For $K=\min(q,m)$ use
\[
 a_k^{(B)}=\frac{\sqrt{(k+1)(m-k)(9r-1)}}m,
 \qquad
 a_k^{(C)}=\frac{\sqrt{(k+1)(m-k)(4r+4)}}m.
\]
For B take $b_{j,k}=B_j(r;\lambda)$ for $k\ge1$. For simple C take
$b_{j,k}=\widetilde B_j(r)$; for refined C take
$b_{j,k}=\mathfrak D_{j,k}(r)$ when $1\le k\le j$ and zero for padding.
For stationary propagation replace $j$ by $q$ in its simple budget.
In every case $b_{j,0}=0$. Start from $w^{(0)}=e_0$ and update
\begin{equation}\label{eq:recurrence}
 w_k^{(j)}=a_{k-1}w_{k-1}^{(j-1)}+
 \frac{b_{j,k}}m w_k^{(j-1)}+a_kw_{k+1}^{(j-1)}.
\end{equation}
Out-of-range terms are omitted. Select $a_k$ from the same chain as the diagonal
budget; do not mix B and C inadvertently. The failure certificate is
\begin{equation}\label{eq:numeric-failure}
 F=\frac r{\varepsilon^{2q}}\sum_{k=0}^K(w_k^{(q)})^2.
\end{equation}
After diagonal costs are precomputed, the recurrence takes $O(qK)$ arithmetic
operations and $O(K)$ storage. Positive interval arithmetic encloses it
without randomized simulation.

\subsection{Computing the exact degree maximum}
For costs $b_d(r)$ from \eqref{c:eq:wick-degree-cost}, let $F_\mathrm{deg}(k,n)$
be the maximum sum for $k$ active labels of exact half-degree $n$. Set
$F_\mathrm{deg}(0,0)=0$ and omit infeasible states. Then
\begin{equation}\label{eq:degree-dp}
 \begin{aligned}
 F_\mathrm{deg}(k,n)
 &=\max_{1\le\ell\le n}
 \{F_\mathrm{deg}(k-1,n-\ell)+b_{2\ell}(r)\},\\
 \mathfrak D_{j,k}(r)&=\max_{k\le n\le2j}F_\mathrm{deg}(k,n).
 \end{aligned}
\end{equation}
For a fixed terminal $q$ only degrees up to $4q$ occur. A direct implementation
costs $O(q^3)$ operations to precompute this finite maximum and $O(q^2)$ storage.
This maximum is over deterministic degree budgets, not over simulated samples.

\subsection{Specified exponent choices}
The following table lists $p_d$ in order $d=2,4,\ldots,4q$. Each choice is fixed
before interval evaluation. The supplementary checks verify
\eqref{c:eq:choice-policy} for each listed degree; claiming a global optimal
exponent is neither needed nor intended.
\begin{table}[H]\centering\small
\setlength{\tabcolsep}{5pt}\renewcommand{\arraystretch}{1.22}
\begin{tabularx}{\textwidth}{@{}rrY@{}}
\toprule
$r$&$q$&$p_2,p_4,\ldots,p_{4q}$\\
\midrule
10&3&4, 6, 9, 11, 13, 16\\
100&5&4, 8, 11, 15, 18, 21, 24, 27, 30, 33\\
1\,000&7&5, 8, 12, 16, 20, 24, 28, 31, 35, 39, 43, 46, 50, 54\\
10\,000&8&5, 8, 12, 16, 20, 24, 28, 32, 36, 40, 44, 48, 52, 56, 60, 64\\
\bottomrule
\end{tabularx}
\caption{Admissible exponent lists used in the refined C-column.}
\label{tab:wick-exponents}
\end{table}
The finite menu $\mathcal P_q$ is only a convenient policy. The theorem does
not require minimization over this menu; a chosen list always gives a valid
certificate, while \eqref{c:eq:choice-policy} supplies the additional comparison
with the simple policy.


\subsection{Verified allowances in the main A/B/C table}
Each entry below is rounded upward to six decimal places at its row count in
Table~\ref{tab:three-main-counts}; all are strictly below $1/2$.
\begin{table}[H]\centering\small
\setlength{\tabcolsep}{12pt}\renewcommand{\arraystretch}{1.2}
\begin{tabular}{@{}rrrr@{}}
\toprule
$r$&\textbf{A}&\textbf{B}&\textbf{C}\\
\midrule
10&0.495576&0.494373&0.499783\\
100&0.497052&0.498858&0.499382\\
1\,000&0.499705&0.499417&0.499353\\
10\,000&0.499744&0.499431&0.498662\\
\bottomrule
\end{tabular}
\caption{Twelve main-table distortion certificates. Computed from proved
inequalities, not estimated from random samples.}\label{tab:allowances}
\end{table}

The companion \texttt{khatri\_rao\_three\_chain\_checks.py} uses 80-digit evaluation
and 70-digit outward-rounded intervals for the twelve principal certificates,
eight supplementary certificates in Appendix~\ref{app:A-wick}, the fixed exponent
policies, and their degree budgets. It records the chain identity explicitly in
the output. All finite exponent choices are fixed before interval validation;
no claim that a selected exponent is globally optimal is needed.
The original verification scripts and outputs are included as archived checks
with their original A/B labels; the new manifest maps them to the present chains.
The analytic arguments, not the finite checks, justify the theorem for arbitrary
subspaces and all admissible parameters.

\subsection{Exact scalar and normalization checks}
For the scalar common-factor model, let $X=g^2h^2-1$. With $(-1)!!=1$,
\begin{equation}\label{eq:scalar-check}
 \mu_k:=\E X^k
 =\sum_{\ell=0}^k\binom{k}{\ell}(-1)^{k-\ell}[(2\ell-1)!!]^2.
\end{equation}
If $M_k^{(j)}=\E(X_1+\cdots+X_j)^k$, independence yields the exact recurrence
\[
 M_k^{(j+1)}=\sum_{\ell=0}^k\binom{k}{\ell}M_{k-\ell}^{(j)}\mu_\ell,
 \qquad M_0^{(0)}=1,\quad M_k^{(0)}=0\ (k>0).
\]
Dividing $M_{2q}^{(m)}$ by $m^{2q}$ gives an exact rational scalar trace moment.
The script checks it against the stage-dependent majorant for
$m\in\{1,2,5\}$ and $1\le q\le6$, giving 18 checks.
At $q=1$, both sides equal $8/m$ exactly because $\E X^2=9-1=8$.
It separately checks
$\norm{J_q^qe_0}^2=(2q-1)!!$ for $1\le q\le8$.

The archived baseline run verifies 16 interval sample certificates, 18 exact
scalar-moment comparisons, eight Jacobi normalizations, eight finite integer
inequalities for the moment scale, and one all-orders rational envelope.
The scalar tests are useful indexing and normalization checks; they are not
proofs for arbitrary subspaces. That generality comes from the analytic
coefficient and operator estimates in the main text.

\subsection{Non-scalar Wick checks and their limits}
The Wick script computes rational polynomial expectations on four coefficient
families: a two-dimensional common-factor subspace, the first two or three
columns of the rational Householder matrix
\[
 H=I_4-\frac{2uu^{\mathsf T}}{30},\qquad u=(1,2,3,4)^{\mathsf T},
\]
reshaped as $2\times2$ matrices, and the full $2\times2$ tensor space.
It checks 24 radial inequalities (orders $2$ through $12$), four exact
covariance Loewner inequalities, and eleven orthogonal mixed-fourth inequalities.
The common-factor radial moments attain \eqref{c:eq:wick-radial} exactly.

For twelve centered fourth-moment comparisons it uses the exact ordered
independent-label expansion
\[
 \E\tr Z^4=
 \frac{\E\tr X^4+(m-1)[2\tr V^2+\E\tr(X_1X_2X_1X_2)]}{m^3},
 \qquad V=\E X^2.
\]
Only the all-equal label and the three pair-pair patterns survive centering;
the displayed alternating word is retained in its original order. This is an
elementary independent-label identity, not a new cumulant theorem. The exact
rational moments are compared numerically with the analytic C-majorant.
These examples check indexing, normalization, and some non-common-factor
geometry; they do not prove a statement for every subspace. That generality
comes from the analytic operator and coefficient arguments.

\clearpage
\section{A supplementary Wick-assisted concentration variant}\label{app:A-wick}
\begin{takeaway}[title={Not a fourth main chain and not the meaning of Chain C}]
The three principal chains are A, B, and C. This appendix separately records
$A^{\mathrm W}$: the same concentration architecture as A supplied with C's
better Gaussian inputs. It preserves the earlier scale-eight result and makes
the constants comparison fair, without silently redefining Chain A.
\end{takeaway}

This variant uses the radial and covariance estimates proved in
Section~\ref{sec:C} but \emph{not} its Hermite propagation. The concentration
step remains Lieb--Tropp; no new matrix concentration theorem is asserted.
Write $\nu_r^{(C)}=4r+5$ for the raw matrix-moment proxy.

\subsection{The Wick-assisted one-row moment estimate}
\begin{lemma}[Wick-refined one-row matrix moments]\label{aw:lem:raw-moments}
For every integer $p\ge2$,
\begin{equation}\label{aw:eq:8-moment}
 \boxed{\E Y^p\preceq\frac{p!}{2}\nu_r^{(C)}[8(r+p)]^{p-2}I_r.}
\end{equation}
For $p\ge3$ the sharper intermediate bound is
\begin{equation}\label{aw:eq:raw-moment}
 \E Y^p\preceq3\sqrt{\kappa_{2p-2}D_{r,2p-2}}\,I_r.
\end{equation}
\end{lemma}
\begin{proof}
\step{1. Retain the rank-one matrix direction.}
For unit $x$, $Y^p=\norm y^{2p-2}yy^{\mathsf T}$ and Cauchy--Schwarz give
\[
 x^{\mathsf T}(\E Y^p)x
 \le(\E\norm y^{4p-4})^{1/2}(\E\ip{x}{y}^4)^{1/2}
 \le3\sqrt{\kappa_{2p-2}D_{r,2p-2}}.
\]
This is \eqref{aw:eq:raw-moment}. The directional factor is bounded by Lemma~\ref{lem:four}.
The case $p=2$ of \eqref{aw:eq:8-moment} is Corollary~\ref{c:cor:wick-variance}.

\step{2. Bound the scalar pairing factor by a factorial sequence.}
Write $H_p=\sqrt{\kappa_{2p-2}}=\sqrt{(4p-5)!!}$. For $p\ge3$,
\[
 H_3=\sqrt{105},\qquad
 \frac{H_{p+1}}{H_p}=\sqrt{(4p-3)(4p-1)}<4p<4(p+1).
\]
Induction yields
\begin{equation}\label{aw:eq:wick-pairingfactor}
 H_p\le\frac{\sqrt{105}}{24}\,p!\,4^{p-2}.
\end{equation}
No finite parameter sweep is being substituted for an all-order argument.

\step{3. Keep the first two radial factors, then use AM--GM.}
There are $2p-4$ remaining factors after $r(r+2)$, and their arithmetic mean is $r+2p-1$. Thus
\[
 \sqrt{D_{r,2p-2}}
 \le\sqrt{r(r+2)}(r+2p-1)^{p-2}
 \le\sqrt{r(r+2)}[2(r+p)]^{p-2}.
\]
Together with \eqref{aw:eq:wick-pairingfactor}, this bounds \eqref{aw:eq:raw-moment} by
\[
 \frac{\sqrt{105}}8\,p!\sqrt{r(r+2)}[8(r+p)]^{p-2}I_r.
\]
Finally $\sqrt{105}/4<3$, $\sqrt{r(r+2)}\le r+1$, and $3(r+1)\le4r+5$. The coefficient is at most $(p!/2)\nu_r^{(C)}$, proving \eqref{aw:eq:8-moment}.
\end{proof}

The matrix direction is retained until expectation; the estimate is on
$\norm{\E Y^p}$, not merely $\E\norm Y^p$. The constant eight is a uniform
factorial envelope, not an asserted optimal scale. Appendix~\ref{app:72} proves the constant-$72$ calculation used by the
unchanged Chain A.

\subsection{Auxiliary truncation and all-order moments}
The untruncated $Y$ need not have a finite positive exponential moment. For example,
when $r=1$ and $y=gh$, $Y=g^2h^2$ has infinite moment-generating function at every
positive argument. Conditional on $g$, the Gaussian integral in $h$ already diverges
whenever $2\theta g^2\ge1$, an event of positive probability.
Thus a positive-matrix exponential theorem cannot be applied to $Y$ directly.

\begin{lemma}[All-order moments of a bounded auxiliary matrix]\label{aw:lem:truncate}
Fix an integer $s\ge2$ and define
\begin{equation}\label{aw:eq:truncation}
 R_s=8(r+s),\qquad L_s=sR_s,\qquad
 W=Y\one_{\{\norm y^2\le L_s\}},\qquad\Sigma=\E W.
\end{equation}
Then
\begin{align}
 \Prob\{W\ne Y\}&\le e^{-2s},\label{aw:eq:clip}\\
 0\preceq I_r-\Sigma&\preceq3e^{-s}I_r,\label{aw:eq:bias}\\
 \E W^p&\preceq\frac{p!}{2}\,\nu_r^{(C)}\,R_s^{p-2}I_r
 \qquad\text{for every integer }p\ge2.\label{aw:eq:all-order}
\end{align}
\end{lemma}
\begin{proof}
\step{1. Bound the event on which the sample is changed.}
At exponent $2s$, Theorem~\ref{c:thm:wick-radial} gives
\[
 \bigl\|\norm y^2\bigr\|_{L_s}
 \le M_s(r)\le s(r+s-1).
\]
Markov's inequality therefore gives
\[
 \Prob\{\norm y^2>L_s\}
 \le\left(\frac{s(r+s-1)}{8s(r+s)}\right)^s
 \le8^{-s}\le e^{-2s}.
\]

\step{2. Charge the covariance bias in each direction.}
For a unit vector $x$, Lemma~\ref{lem:four} and Cauchy--Schwarz imply
\[
 x^{\mathsf T}(I_r-\Sigma)x
 =\E\bigl[\ip{x}{y}^2\one_{\{\norm y^2>L_s\}}\bigr]
 \le3\Prob\{\norm y^2>L_s\}^{1/2}\le3e^{-s}.
\]
The integrand is nonnegative, proving both sides of \eqref{aw:eq:bias}.

\step{3. Prove the matrix moment bound through order $s$.}
For $p=2$, $\E W^2\preceq\E Y^2\preceq\nu_r^{(C)}I_r$.
For $3\le p\le s$, the scalar event in \eqref{aw:eq:truncation} gives
$W^p\preceq Y^p$ pointwise. Lemma~\ref{aw:lem:raw-moments} and $p\le s$ yield
\[
 \E W^p\preceq\frac{p!}{2}\,\nu_r^{(C)}\,[8(r+p)]^{p-2}I_r
 \preceq\frac{p!}{2}\,\nu_r^{(C)}\,R_s^{p-2}I_r.
\]

\step{4. Extend to every order required by the exponential series.}
For an integer $p>s$, spectral calculus on $0\preceq W\preceq L_sI_r$ gives
\[
 W^p\preceq L_s^{p-s}W^s.
\]
Consequently,
\[
 \E W^p\preceq\frac{s!}{2}\,\nu_r^{(C)}\,R_s^{s-2}L_s^{p-s}I_r
 =\frac{s!s^{p-s}}{2}\,\nu_r^{(C)}\,R_s^{p-2}I_r
 \preceq\frac{p!}{2}\,\nu_r^{(C)}\,R_s^{p-2}I_r.
\]
The last inequality follows by comparing each factor of
$(s+1)(s+2)\cdots p$ with $s$. This proves all of \eqref{aw:eq:all-order},
not just the first $s$ moments.
\end{proof}
\begin{takeaway}
The almost-sure cutoff height is $L_s=sR_s$, but the matrix moment-growth scale is
only $R_s$. Charging $L_s$ as the bounded-summand scale would lose the factor $s$.
\end{takeaway}


\subsection{Restore the original samples and compare with C}
Apply the concentration lemma already proved in Chain A,
Lemma~\ref{lem:bernstein}, with $v=4r+5$ and $R=8(r+s)$.
Choose $t=\log(4r/\delta)$ and $s=\lceil\log(4m/\delta)\rceil$.
The all-order estimate above supplies its hypotheses. A union bound for
clipping gives at most $me^{-2s}\le\delta/16$; the concentration event costs
$2re^{-t}=\delta/2$. The covariance bias is at most $3e^{-s}$.
On the good event the original samples equal the auxiliary ones, so
\begin{equation}\label{aw:eq:tail}
 \mathbb P\left\{\|Z\|>
 U_{A^{\mathrm W}}:=
 \sqrt{\frac{2(4r+5)t}{m}}+
 \frac{16(r+s)t}{m}+3e^{-s}\right\}\le\delta.
\end{equation}
The sample-order absorption proved in Proposition~\ref{prop:A-size} applies
with these smaller constants. This is still the original Gaussian sketch,
not a clipped estimator.

\begin{table}[H]\centering\small
\setlength{\tabcolsep}{10pt}\renewcommand{\arraystretch}{1.2}
\begin{tabular}{@{}rrrr@{}}
\toprule
$r$&\textbf{Supplementary $A^{\mathrm W}$}&\textbf{C simple}&\textbf{C refined}\\
\midrule
10&13\,200&19\,700&12\,800\\
100&98\,100&148\,000&111\,000\\
1\,000&1\,100\,000&1\,620\,000&1\,250\,000\\
10\,000&12\,800\,000&19\,300\,000&15\,000\,000\\
\bottomrule
\end{tabular}
\caption{A fair supplementary comparison after giving the concentration route
the Wick input as well. The refined C-column repeats the primary C certificate;
A in the main table continues to mean the original concentration baseline.}
\label{tab:AWC}
\end{table}
The corresponding upward-rounded allowances for $A^{\mathrm W}$ are
$0.499192$, $0.499736$, $0.497858$, $0.498718$; those for simple C are
$0.499266$, $0.498479$, $0.499814$, $0.499373$.
The same $q$ values as in Table~\ref{tab:three-main-counts} are used in both C-columns.
All entries are sufficient certificates, not global minima. The concentration
variant is smaller than refined C in three rows, while refined C is smaller
at $r=10$. The large improvement of C over the original A/B baselines is
therefore not evidence of universal superiority over concentration.

\begin{thebibliography}{99}
\raggedright
\addcontentsline{toc}{section}{References}
\bibitem{Framework}
D. Heidary, \emph{Graded Multiplication Operators for Random-Matrix Moments},
Version 11, exposition and attribution revision, 16 September 2026.
Unpublished supplied manuscript \texttt{v11\_attribution\_revised.tex}.
Organizing reference for the domain, coefficient, vacuum-majorant, and energy
interfaces; application-specific steps are proved in this note.

\bibitem{Sun}
Y. Sun, Y. Guo, J. A. Tropp, and M. Udell,
\emph{Tensor random projection for low-memory dimension reduction},
NeurIPS Workshop on Relational Representation Learning, 2018.
Later arXiv posting: \href{https://arxiv.org/abs/2105.00105v1}{arXiv:2105.00105v1},
30 April 2021. This is a workshop paper, not a main-conference NeurIPS proceedings article.

\bibitem{Ahle}
T. D. Ahle, M. Kapralov, J. B. T. Knudsen, R. Pagh, A. Velingker,
D. P. Woodruff, and A. Zandieh,
\emph{Oblivious sketching of high-degree polynomial kernels},
in \emph{Proceedings of the 2020 ACM--SIAM Symposium on Discrete Algorithms
(SODA)}, 2020, 141--160.
\href{https://doi.org/10.1137/1.9781611975994.9}{DOI: 10.1137/1.9781611975994.9}.
\href{https://arxiv.org/abs/1909.01410}{arXiv:1909.01410}.

\bibitem{BK}
Z. Bujanovi\'c and D. Kressner,
\emph{Norm and trace estimation with random rank-one vectors},
SIAM Journal on Matrix Analysis and Applications \textbf{42}(1) (2021), 202--223.
\href{https://doi.org/10.1137/20M1331718}{DOI: 10.1137/20M1331718}.
\href{https://arxiv.org/abs/2004.06433v2}{arXiv:2004.06433v2}.

\bibitem{BGKL}
Z. Bujanovi\'c, L. Grubi\v{s}i\'c, D. Kressner, and H. Y. Lam,
\emph{Subspace embedding with random Khatri--Rao products and its application to
eigensolvers}, IMA Journal of Numerical Analysis \textbf{46}(3) (2026), 1328--1355.
First published online 26 June 2025.
\href{https://doi.org/10.1093/imanum/draf043}{DOI: 10.1093/imanum/draf043}.
\href{https://arxiv.org/abs/2405.11962v1}{arXiv:2405.11962v1} (2024).

\bibitem{SVBv1}
A. K. Saibaba, B. D. Verma, and G. Ballard,
\emph{Improved analysis of Khatri--Rao random projections and applications},
\href{https://arxiv.org/abs/2507.23207v1}{arXiv:2507.23207v1}, 31 July 2025.
Theorem 11 and Appendix B are cited as the historical OSE claim and proof attempt,
not as a theorem used in this note; see~\cite{SVBv2,BM} and the discussion in
Section~\ref{sec:attribution-intro}.

\bibitem{SVBv2}
A. K. Saibaba, B. D. Verma, and G. Ballard,
\emph{Improved analysis of Khatri--Rao random projections and applications},
\href{https://arxiv.org/abs/2507.23207v2}{arXiv:2507.23207v2}, 1 April 2026.
Revised version; the earlier Theorem 11 OSE assertion is omitted.

\bibitem{BM}
L. Beretta and C. Musco,
\emph{A tight analysis of Khatri--Rao oblivious subspace embeddings},
\href{https://arxiv.org/abs/2608.28094v2}{arXiv:2608.28094v2}, 8 September 2026.
Theorem 1.3, Section 2.3, and acknowledgments.
The rate comparison is restricted to the real-Gaussian order-two specialization.

\bibitem{Nelson}
E. Nelson, \emph{The free Markoff field},
Journal of Functional Analysis \textbf{12}(2) (1973), 211--227.
\href{https://doi.org/10.1016/0022-1236(73)90025-6}{DOI: 10.1016/0022-1236(73)90025-6}.

\bibitem{Gross}
L. Gross, \emph{Logarithmic Sobolev inequalities},
American Journal of Mathematics \textbf{97}(4) (1975), 1061--1083.
\href{https://doi.org/10.2307/2373688}{DOI: 10.2307/2373688}.

\bibitem{GrossSurvey}
L. Gross, \emph{Logarithmic Sobolev inequalities and contractivity properties of
semigroups}, in G. Dell'Antonio and U. Mosco (eds.), \emph{Dirichlet Forms},
Lecture Notes in Mathematics \textbf{1563}, Springer, 1993, 54--88.
\href{https://doi.org/10.1007/BFb0074091}{DOI: 10.1007/BFb0074091}.

\bibitem{AidaStroock}
S. Aida and D. Stroock,
\emph{Moment estimates derived from Poincar\'e and logarithmic Sobolev inequalities},
Mathematical Research Letters \textbf{1}(1) (1994), 75--86.
\href{https://doi.org/10.4310/MRL.1994.v1.n1.a9}{DOI: 10.4310/MRL.1994.v1.n1.a9}.

\bibitem{ALM}
R. Adamczak, R. Lata\l a, and R. Meller,
\emph{Hanson--Wright inequality in Banach spaces},
Annales de l'Institut Henri Poincar\'e, Probabilit\'es et Statistiques
\textbf{56}(4) (2020), 2356--2376.
\href{https://doi.org/10.1214/19-AIHP1041}{DOI: 10.1214/19-AIHP1041}.
\href{https://arxiv.org/abs/1811.00353v3}{arXiv:1811.00353v3}, Corollary 12.

\bibitem{Isserlis}
L. Isserlis,
\emph{On a formula for the product-moment coefficient of any order of a normal
frequency distribution in any number of variables},
Biometrika \textbf{12}(1--2) (1918), 134--139.
\href{https://doi.org/10.1093/biomet/12.1-2.134}{DOI: 10.1093/biomet/12.1-2.134}.

\bibitem{Lieb}
E. H. Lieb, \emph{Convex trace functions and the Wigner--Yanase--Dyson conjecture},
Advances in Mathematics \textbf{11}(3) (1973), 267--288, Theorem 6.
\href{https://doi.org/10.1016/0001-8708(73)90011-X}{DOI: 10.1016/0001-8708(73)90011-X}.
The trace-concavity consequence is stated in \eqref{eq:lieb}; see also
\cite[Theorem~3.2]{Tropp}.

\bibitem{Tropp}
J. A. Tropp, \emph{User-friendly tail bounds for sums of random matrices},
Foundations of Computational Mathematics \textbf{12} (2012), 389--434.
\href{https://doi.org/10.1007/s10208-011-9099-z}{DOI: 10.1007/s10208-011-9099-z}.
\href{https://arxiv.org/abs/1004.4389v7}{arXiv:1004.4389v7}.
Proposition 3.1, Theorem 3.2, Corollary 3.3, Lemma 3.4, Theorem 6.2,
and Lemma 6.8 describe the Laplace and subexponential-moment mechanisms used here.

\bibitem{Wiener}
N. Wiener, \emph{The homogeneous chaos},
American Journal of Mathematics \textbf{60}(4) (1938), 897--936.
\href{https://doi.org/10.2307/2371268}{DOI: 10.2307/2371268}.

\bibitem{AccardiBozejko}
L. Accardi and M. Bo\.{z}ejko,
\emph{Interacting Fock spaces and Gaussianization of probability measures},
Infinite Dimensional Analysis, Quantum Probability and Related Topics
\textbf{1}(4) (1998), 663--670.
\href{https://doi.org/10.1142/S0219025798000363}{DOI: 10.1142/S0219025798000363}.

\bibitem{GolubWelsch}
G. H. Golub and J. H. Welsch,
\emph{Calculation of Gauss quadrature rules},
Mathematics of Computation \textbf{23}(106) (1969), 221--230.
\href{https://doi.org/10.1090/S0025-5718-69-99647-1}{DOI: 10.1090/S0025-5718-69-99647-1}.
An earlier report is Stanford Computer Science Technical Report CS 81 (1967).

\bibitem{BLNvH}
A. S. Bandeira, K. Lucca, P. Nizi\'c-Nikolac, and R. van Handel,
\emph{Matrix chaos inequalities and chaos of combinatorial type},
\href{https://arxiv.org/abs/2412.18468v2}{arXiv:2412.18468v2}, 31 March 2025.
Sections 2--3 give the coefficient-flattening framework; Section 4.1 treats
Khatri--Rao full-matrix norm bounds. These are methodological context, not a
substitute for the arbitrary-fixed-subspace covariance proof in this note.
\bibitem{Janson}
S. Janson, \emph{Gaussian Hilbert Spaces}, Cambridge Tracts in Mathematics
\textbf{129}, Cambridge University Press, 1997.
\href{https://doi.org/10.1017/CBO9780511526169}{DOI: 10.1017/CBO9780511526169}.
Chapters 2--5: classical Gaussian chaos, Wick structure, and hypercontractivity.

\bibitem{RMPS}
\emph{Gaussian rMPS Moments: Two Self-Contained Proof Chains}.
Supplied attributed project manuscript,
\texttt{rmps\_two\_proofs\_attributed.tex}, 16 September 2026.
Role: the scalar rank-one Gaussian envelope and the projected radial-moment
mechanism; not an external theorem about the present Khatri--Rao model.
Its public Gaussian and tensor-train attributions are not replaced by this
project citation.

\bibitem{WickNote}
\emph{Wick refinements of Gaussian Khatri--Rao embedding constants}.
Supplied research supplement, \texttt{khatri\_rao\_wick\_refinement.tex},
16 September 2026. The radial, mixed-fourth, scale-eight, and refined-degree
calculations are integrated and proved in this document. This is internal
provenance, not evidence of external review or publication priority.

\end{thebibliography}
\end{document}
